\chapter{Complexity — Query Reshape}
%%% generated by the blueprint pipeline from TCSlib.Complexity.SpaceComplexity.Logspace.QueryReshape
\section{Overview}

Declarations of the Lean module \texttt{TCSlib.Complexity.SpaceComplexity.Logspace.QueryReshape}.

\section{Declarations}

\begin{definition}[Removing an encoded index]
\lean{Complexity.dropIdx}
\label{Complexity.dropIdx}
\leanok
\uses{Turing.pairDecode, Turing.pairEncode}
Define a function on finite Boolean strings as follows: decode the string as a pair
$(u,b')$, then decode $u$ as a pair $(w,b)$. If both decodings succeed, return the
encoding of $(w,b')$; if either fails, return the empty string.
\end{definition}

\begin{theorem}[Dropping an index from a nested encoding]
\lean{Complexity.dropIdx_pairEncode}
\label{Complexity.dropIdx_pairEncode}
\leanok
\uses{Complexity.dropIdx, Turing.pairDecode_pairEncode, Turing.pairEncode}
\difficulty{1}
For any finite bit strings $w,b,b'$, applying the index-dropping operation to the nested
encoding of $w,b,b'$ returns the pair encoding of $w$ and $b'$.
\end{theorem}

\begin{theorem}[Length bound for dropping an index]
\lean{Complexity.dropIdx_length_le}
\label{Complexity.dropIdx_length_le}
\leanok
\uses{Complexity.dropIdx, Turing.eq_pairEncode_of_pairDecode, Turing.length_pairEncode}
\difficulty{4}
For every finite bit string $z$, applying the operation that parses the nested pair
encoding and drops the inner second component, returning the empty string when parsing
fails, produces a bit string whose length is at most the length of $z$.
\end{theorem}

\begin{definition}[Composing a function with encoded data]
\lean{Complexity.compQ}
\label{Complexity.compQ}
\leanok
\uses{Complexity.LogProg.bitsVal, Turing.pairDecode, Turing.pairEncode}
Let $f$ be a function from finite Boolean strings to finite Boolean strings. Define its
encoded composition on a finite Boolean string $w$ by decoding $w$ as $(x,b)$; if the
decoding succeeds and the binary expansion of the natural number encoded by $b$ equals
$b$, return the encoding of $(f(x),b)$, and otherwise return the empty string.
\end{definition}

\begin{theorem}[Transformation on encoded queries]
\lean{Complexity.compQ_pairEncode}
\label{Complexity.compQ_pairEncode}
\leanok
\uses{Complexity.LogProg.bitsVal_bits, Complexity.compQ, Turing.pairDecode_pairEncode, Turing.pairEncode}
\difficulty{1}
Let $f$ map bit strings to bit strings, let $x$ be a bit string, and let $i$ be a
natural number. Applying the transformation to the encoding of
$(x,\operatorname{bin}(i))$ returns the encoding of $(f(x),\operatorname{bin}(i))$,
where $\operatorname{bin}(i)$ is the binary expansion of $i$.
\end{theorem}

\begin{theorem}[Membership under index-language composition]
\lean{Complexity.mem_indexLang_compQ_iff}
\label{Complexity.mem_indexLang_compQ_iff}
\leanok
\uses{Complexity.compQ, Complexity.compQ_pairEncode, Complexity.indexLang, Complexity.pairEncode_ne_nil, Turing.eq_pairEncode_of_pairDecode, Turing.pairEncode_injective}
\difficulty{5}
Let $f$ map binary strings to binary strings, let $p$ be a predicate on a binary string
and a natural number, and let $w$ be a binary string. Then $w$ is an encoding of a pair
$(x,\operatorname{bits}(i))$ with $p(f(x),i)$ true if and only if applying the following
transformation to $w$ produces a string that is an encoding of some pair
$(y,\operatorname{bits}(j))$ with $p(y,j)$ true: parse $w$ as a pair $(x,b)$; if parsing
succeeds and $b$ is the canonical bit expansion of the natural number represented by
$b$, output the encoding of $(f(x),b)$, and otherwise output the empty string.
\end{theorem}

\begin{theorem}[Polynomial output-length bound]
\lean{Complexity.compQ_length_le}
\label{Complexity.compQ_length_le}
\leanok
\uses{Complexity.compQ, Turing.eq_pairEncode_of_pairDecode, Turing.length_pairEncode}
\difficulty{5}
Let $f$ map finite bit strings to finite bit strings, and let $C,c$ be natural numbers.
If every bit string $x$ satisfies $|f(x)|\le C(|x|+1)^c$, then for every bit string $w$,
the output of the following partial decoding and transformation has length at most
$(2C+3)(|w|+1)^{c+1}$: parse $w$ as a pair $(x,b)$; if parsing succeeds and $b$ is the
binary expansion of its value, return the encoding of $(f(x),b)$, and otherwise return
the empty string.
\end{theorem}

\begin{lemma}[Prefix address under index dropping]
\lean{Complexity.dropIdx_prefix_address}
\label{Complexity.dropIdx_prefix_address}
\leanok
\uses{Complexity.dropIdx, Complexity.dropIdx_pairEncode, Complexity.pairEncode_getD, Turing.length_pairEncode, Turing.pairEncode}
\difficulty{4}
Let $w,b,b'$ and $\mathrm{tail}$ be finite bit strings, and let $j$ be a natural number
with $j<2|w|+2$. The bit at index $j$ of the result of dropping the index from the
encoding of the encoding of $w,b$ paired with $b'$ equals the bit at index $4j$ of the
encoding of that nested string paired with $\mathrm{tail}$. Out-of-range lookups are
interpreted as the default bit $0$.
\end{lemma}

\begin{lemma}[Suffix address under index dropping]
\lean{Complexity.dropIdx_suffix_address}
\label{Complexity.dropIdx_suffix_address}
\leanok
\uses{Complexity.dropIdx, Complexity.dropIdx_pairEncode, Complexity.pairEncode_getD, Turing.length_pairEncode, Turing.pairEncode}
\difficulty{4}
Let $w,b,b',\mathrm{tail}$ be finite bit strings, and let $j$ be a natural number with
$j<|b'|$. The bit at index $2|w|+2+j$ of the result of dropping the index from the
encoding of the encoding of $w,b$ paired with $b'$ equals the bit at index $4|\langle
w,b\rangle|+4+2j$ of the encoding of that nested string paired with $\mathrm{tail}$,
where $\langle w,b\rangle$ denotes the encoding of $w$ and $b$. Out-of-range lookups are
interpreted as the default bit $0$.
\end{lemma}

\begin{lemma}[Bound from binary length]
\lean{Complexity.query_index_bound}
\label{Complexity.query_index_bound}
\leanok
\difficulty{4}
For natural numbers $r$ and $n$, if the binary expansion of $r$ has length at most the
binary expansion of $n$, then $r \le 2n$.
\end{lemma}

\begin{lemma}[Value of a binary suffix]
\lean{Complexity.query_suffix_value_le}
\label{Complexity.query_suffix_value_le}
\leanok
\uses{Complexity.LogProg.bitsVal}
\difficulty{4}
For every finite Boolean word $w$ and every $j\in\mathbb{N}$, the value of the suffix
obtained by deleting the first $j$ bits of $w$ is at most the value of $w$, where the
value is computed with the first bit as the least significant bit.
\end{lemma}

\begin{lemma}[Bound on the Value of a Canonical Suffix]
\lean{Complexity.query_accumulator_bound}
\label{Complexity.query_accumulator_bound}
\leanok
\uses{Complexity.LogProg.Canon, Complexity.LogProg.bitsVal, Complexity.LogProg.canon_eq_bits, Complexity.query_index_bound, Complexity.query_suffix_value_le}
\difficulty{3}
Let $w$ be a Boolean word, let $n,j\in\mathbb{N}$, and suppose that $w$ is empty or its
last bit is true. If the length of $w$ is at most the binary bit length of $n$, then the
value represented by the suffix of $w$ after dropping its first $j$ bits is at most
$2n$.
\end{lemma}

\begin{lemma}[Overflow in a dropped-index query]
\lean{Complexity.dropIdx_query_overflow}
\label{Complexity.dropIdx_query_overflow}
\leanok
\uses{Complexity.LogProg.length_bits_mono, Complexity.dropIdx, Complexity.dropIdx_length_le, Turing.length_pairEncode, Turing.pairEncode}
\difficulty{3}
Let $z$ be a finite bit string and $r$ a natural number. Suppose the binary
representation of the length of the encoding of $z$ paired with the binary
representation of $r$ has fewer bits than the binary representation of $r$. Then the
result of dropping the index from $z$ has default bit $0$ at index $r$, and $r$ is not
less than the length of that result.
\end{lemma}

\begin{definition}[Drop machine states]
\lean{Complexity.DropState}
\label{Complexity.DropState}
\leanok
A drop-machine state is one of the following: a length test or length increment; a copy
state with a stage in $\{0,\ldots,7\}$ and a Boolean flag; a scan or first state with a
layer in $\{0,1,2\}$; a second or second-read state with such a layer and a Boolean
flag; a scan-increment state with a layer and two indices in $\{0,1,2,3\}$; a tail-test
state with a Boolean flag, or one of the tail-check, tail-read, or tail-done states; a
tail-half or tail-increment state with a Boolean flag; a pointer-clearing state with a
stage in $\{0,1,2,3,4\}$; a $z$-decrement state with a step in $\{0,1\}$, or a
$u$-decrement state with a step in $\{0,1,2,3\}$; one of the decode, decode-down, or
decode-read states; a decode-copy, decode-copy-increment, decode-clear, decode-scale,
decode-take, or decode-add state with a Boolean flag; a decode-increment state with two
Boolean flags; a seek, match-index, or next-index state with a Boolean flag; a walk
state with a Boolean flag and a step in $\{0,1,2,3\}$; a hit state; or an answer state
with a Boolean flag.
\end{definition}

\begin{definition}[Drop copy transition table]
\lean{Complexity.dropCopy}
\label{Complexity.dropCopy}
\leanok
\uses{Complexity.DropState}
Define a transition table on indices $0,1,\ldots,7$ whose entries are triples consisting
of two indices in $\{0,\ldots,8\}$ and a drop-machine state. In order, the entries are
$(8,0,\text{copy stage }1\text{ with flag false})$, $(7,1,\text{tail test with flag
true})$, $(0,1,\text{decode})$, $(2,8,\text{$z$-decrement step }0)$, $(3,0,\text{copy
stage }5\text{ with flag false})$, $(4,0,\text{$u$-decrement step }0)$,
$(5,0,\text{clear pointer stage }3)$, and $(0,3,\text{seek with flag true})$.
\end{definition}

\begin{definition}[Drop layer stride]
\lean{Complexity.dropStride}
\label{Complexity.dropStride}
\leanok
For a layer indexed by $0,1,$ or $2$, define its stride to be $1,2,$ or $4$,
respectively.
\end{definition}

\begin{definition}[Drop layer limit]
\lean{Complexity.dropLimit}
\label{Complexity.dropLimit}
\leanok
For a layer indexed by $0,1,$ or $2$, define its limit to be $1,2,$ or $4$,
respectively, regarded as an element of $\{0,1,\ldots,8\}$.
\end{definition}

\begin{definition}[Drop scan completion state]
\lean{Complexity.dropScanDone}
\label{Complexity.dropScanDone}
\leanok
\uses{Complexity.DropState}
For a layer indexed by $0,1,$ or $2$, define the scan-completion state to be the copy
state with flag false and stage $0,4,$ or $6$, respectively.
\end{definition}

\begin{definition}[Drop scan advance]
\lean{Complexity.dropAdvance}
\label{Complexity.dropAdvance}
\leanok
\uses{Complexity.DropState, Complexity.dropScanDone, Complexity.dropStride}
Given a layer $\ell\in\{0,1,2\}$, a choice $c\in\{0,1,2,3\}$, and a natural number $i$,
define the next state as the scan-increment state for layer $\ell$, choice $c$, and step
$i$ if $i$ is less than the stride of $\ell$. Otherwise, according as $c=0,1,2,3$,
define it to be the second state for layer $\ell$ with flag false, the second state with
flag true, the scan state for layer $\ell$, or the scan-completion state for layer
$\ell$.
\end{definition}

\begin{definition}[Drop tail advance]
\lean{Complexity.dropTailAdvance}
\label{Complexity.dropTailAdvance}
\leanok
\uses{Complexity.DropState}
Given a Boolean flag $b$ and a natural number $i$, define the next state to be the
tail-increment state with flag $b$ if $i=0$, and the tail-test state with flag $b$
otherwise.
\end{definition}

\begin{definition}[Drop decode advance]
\lean{Complexity.dropDecodeAdvance}
\label{Complexity.dropDecodeAdvance}
\leanok
\uses{Complexity.DropState}
Given a Boolean flag $b$ and a natural number $i$, define the next state to be the
decode-increment state with flags $b$ and false if $i=0$, the decode-increment state
with flags $b$ and true if $i=1$, and the decode-scale state with flag $b$ otherwise.
\end{definition}

\begin{definition}[Bounded Drop-Walk State]
\lean{Complexity.dropWalk}
\label{Complexity.dropWalk}
\leanok
\uses{Complexity.DropState}
For a Boolean flag $s$ and a natural number $i$, define a state that is the $i$th walk
state with flag $s$ when $i<2$ if $s$ is true or $i<4$ if $s$ is false; otherwise define
it to be the seek state with flag $s$.
\end{definition}

\begin{definition}[Drop Oracle Machine Transition Function]
\lean{Complexity.dropARM}
\label{Complexity.dropARM}
\leanok
\uses{Complexity.DropState, Complexity.LogProg.ARM, Complexity.dropAdvance, Complexity.dropCopy, Complexity.dropLimit, Complexity.dropWalk}
For each Boolean parameter $b$, define a transition function for a machine with nine
registers, two oracle languages, and the finite control states listed in the
declaration. At each control state, the function prescribes the corresponding register
operation, branch, oracle call, or halt exactly as specified by the transition rules;
when the search reaches a match, it returns true if $b$ is true and otherwise queries
the first oracle before choosing its answer.
\end{definition}

\begin{definition}[Drop Oracle Language Pair]
\lean{Complexity.dropBank}
\label{Complexity.dropBank}
\leanok
\uses{Complexity.inputBitLang, Complexity.inputLenLang}
Define a pair of Boolean languages indexed by $j\in\{0,1\}$: the language at index $0$
is the input-bit language, and the language at index $1$ is the input-length language.
\end{definition}

\begin{definition}[Drop Index Language]
\lean{Complexity.dropLang}
\label{Complexity.dropLang}
\leanok
\uses{Complexity.dropIdx, Complexity.indexLang}
For each Boolean length-mode flag, define a language by accepting a pair consisting of
an input $z$ and a query index $q$ exactly when, in length mode, $q$ is less than the
length of the index list associated to $z$, or, otherwise, the entry at index $q$ in
that list (defaulting to false when absent) is true.
\end{definition}

\begin{lemma}[Preconditions for Drop Instructions]
\lean{Complexity.drop_preS}
\label{Complexity.drop_preS}
\leanok
\uses{Complexity.DropState, Complexity.LogProg.AConf, Complexity.LogProg.PreS, Complexity.dropARM}
\difficulty{5}
For every mode, input bit string $h$, and configuration $c$ of the specified
nine-register computation, the following instruction-dependent conditions hold: if the
active instruction compares two registers, their indices are distinct; if it compares a
register’s binary representation with the payload of a plain encoded pair, the input has
that plain-pair form with a canonical payload; if it compares against either component
of a nested encoded pair, the input has that nested-pair form with both payloads
canonical; and if it calls an oracle using a list of registers, the list has no
repetitions.
\end{lemma}

\begin{definition}[Drop Oracle Indicator]
\lean{Complexity.dropOracle}
\label{Complexity.dropOracle}
\leanok
\uses{Complexity.dropBank, Turing.MultiTapeTM.indicator}
For an oracle index $j\in\{0,1\}$ and a Boolean word $w$, define the oracle answer to be
the indicator of whether $w$ belongs to the language at index $j$ in the drop oracle
pair.
\end{definition}

\begin{definition}[Good Drop Configuration]
\lean{Complexity.DropGood}
\label{Complexity.DropGood}
\leanok
\uses{Complexity.DropState, Complexity.LogProg.AConf, Complexity.LogProg.PreS, Complexity.dropARM}
For a Boolean mode, Boolean history $h$, and configuration $a$ of the nine-register
machine, call the configuration good when it is a predecessor configuration for that
machine after history $h$, and every register value is at most $4|h|+4$.
\end{definition}

\begin{definition}[Reachable Drop Configurations]
\lean{Complexity.DropReach}
\label{Complexity.DropReach}
\leanok
\uses{Complexity.DropGood, Complexity.LogProg.AReach, Complexity.dropARM, Complexity.dropOracle}
For a Boolean mode and Boolean history $h$, define the reachable configurations to be
those reachable for the drop machine using the drop oracle, subject to the
good-configuration condition for that mode and history.
\end{definition}

\begin{definition}[Halting Drop Configurations]
\lean{Complexity.DropHalt}
\label{Complexity.DropHalt}
\leanok
\uses{Complexity.DropGood, Complexity.LogProg.AHalt, Complexity.dropARM, Complexity.dropOracle}
For a Boolean mode and Boolean history $h$, define the halting configurations to be the
reachable configurations for the drop machine using the drop oracle that satisfy the
good-configuration condition for that mode and history and are halting configurations.
\end{definition}

\begin{definition}[Register values for dropping]
\lean{Complexity.dropValues}
\label{Complexity.dropValues}
\leanok
For a natural number $n$ and six natural-number parameters $\mathrm{ptr}$,
$\mathrm{zend}$, $\mathrm{retained}$, $\mathrm{inner}$, $\mathrm{kept}$,
$\mathrm{query}$, $\mathrm{temp}$, and $\mathrm{tail}$, define a function on the
nine-element set $\{0,1,\ldots,8\}$ by assigning, in order, the values $\mathrm{ptr}, n,
\mathrm{zend}, \mathrm{retained}, \mathrm{inner}, \mathrm{kept}, \mathrm{query},
\mathrm{temp}, \mathrm{tail}$.
\end{definition}

\begin{lemma}[Oracle returns an input bit]
\lean{Complexity.drop_bit}
\label{Complexity.drop_bit}
\leanok
\uses{Complexity.dropOracle, Complexity.indicator_inputBitLang, Complexity.inputBitLang, Turing.MultiTapeTM.indicator, Turing.pairEncode}
\difficulty{1}
For every bit string $h$ and natural number $i$, the bit oracle, queried on the encoding
of $h$ paired with the binary representation of $i$, returns the $i$th bit of $h$, or
false if $i$ is beyond its length.
\end{lemma}

\begin{lemma}[Length Oracle on Encoded Input]
\lean{Complexity.drop_len}
\label{Complexity.drop_len}
\leanok
\uses{Complexity.dropOracle, Complexity.indicator_inputLenLang, Complexity.inputLenLang, Turing.MultiTapeTM.indicator, Turing.pairEncode}
\difficulty{1}
For every bit string $h$ and natural number $i$, the length oracle, evaluated on
$\langle h,\mathrm{bits}(i)\rangle$, returns true exactly when $i<|h|$.
\end{lemma}

\begin{lemma}[Halting from an Answer State]
\lean{Complexity.drop_ret}
\label{Complexity.drop_ret}
\leanok
\uses{Complexity.DropHalt, Complexity.LogProg.AHalt.ret, Complexity.drop_preS}
\difficulty{1}
Let $m$ and $b$ be Boolean values, let $h$ be a bit string, and let $v$ assign a natural
number to each of nine registers, with every register value at most $4|h|+4$. Starting
from the answer state for $b$, with register values $v$ and no answer yet recorded, the
computation halts with answer $b$, and every configuration before halting satisfies the
stated computation invariant.
\end{lemma}

\begin{lemma}[Halting from the Hit State]
\lean{Complexity.drop_hit}
\label{Complexity.drop_hit}
\leanok
\uses{Complexity.DropHalt, Complexity.LogProg.AHalt.ret, Complexity.LogProg.AHalt.step, Complexity.LogProg.astep, Complexity.LogProg.vword_whole₁, Complexity.dropARM, Complexity.drop_bit, Complexity.drop_preS, Complexity.drop_ret}
\difficulty{4}
Let $h$ be a bit string, let $mode$ be a Boolean mode, and let $v$ assign natural
numbers to the nine registers, each at most $4|h|+4$. Starting in state $hit$ with
register values $v$ and no answer, the machine in mode $mode$ halts with answer true if
$mode$ is true, and otherwise with the bit of $h$ at position $v(0)$, interpreted as
false when that position is outside $h$. Every configuration before halting satisfies
the machine's safety condition and the same register bound.
\end{lemma}

\begin{lemma}[A bounded copy-register run]
\lean{Complexity.drop_copy_run}
\label{Complexity.drop_copy_run}
\leanok
\uses{Complexity.DropReach, Complexity.LogProg.registerCopy_loop, Complexity.dropARM, Complexity.dropCopy, Complexity.dropOracle}
\difficulty{1}
Fix a mode, a Boolean input $h$, a stage among eight copy stages, and an initial
assignment $v$ of natural numbers to nine registers. Suppose every register initially
holds at most $4|h|+4$, and the value in the stage's source register is at most the
value in its destination register. Then the specified register machine, with its
associated oracle, has a run from the configuration at that stage with no answer yet to
the configuration at the next stage with the source register replaced by the destination
register's initial value and no answer yet; every configuration before the final one
satisfies the machine's validity conditions and keeps every register at most $4|h|+4$.
\end{lemma}

\begin{lemma}[Pair-Scan Run]
\lean{Complexity.drop_pair_run}
\label{Complexity.drop_pair_run}
\leanok
\uses{Complexity.DropReach, Complexity.LogProg.alignedPair_loop, Complexity.dropARM, Complexity.dropAdvance, Complexity.dropLimit, Complexity.dropOracle, Complexity.dropScanDone, Complexity.dropStride, Complexity.drop_bit, Turing.pairDecode}
\difficulty{4}
Let $m\in\{\mathrm{false},\mathrm{true}\}$, let $h$ and $w$ be bit strings, let
$\ell\in\{0,1,2\}$, and let $v$ assign a natural number to each of nine registers.
Suppose every register value is at most $4|h|+4$, the limit register equals the pointer
register plus $s_\ell|w|$, where $s_0=1$, $s_1=2$, and $s_2=4$, and, for every $i<|w|$,
the bit of $h$ at position $v(0)+s_\ell i$ (taken as $0$ when out of range) equals the
bit of $w$ at position $i$ (also taken as $0$ when out of range). If parsing $w$ as an
aligned pair fails, the computation from the scan state reaches a false-answer state
after advancing the pointer by $s_\ell t$ for some $t\le |w|$. If the parse succeeds
with first component $a$, it reaches the scan-complete state with the pointer advanced
by $s_\ell(2|a|+2)$. In either case, the computation invariant holds at every step
before the reached state.
\end{lemma}

\begin{lemma}[Evaluation of the drop-value vector]
\lean{Complexity.dropValues_apply}
\label{Complexity.dropValues_apply}
\leanok
\uses{Complexity.dropValues}
\difficulty{1}
Let $n,p,z,s,u,f,q,t,a$ be natural numbers, and form the length-nine vector
$(p,n,z,s,u,f,q,t,a)$. For every index $r$ from $0$ to $8$, evaluating the vector at $r$
gives the corresponding entry.
\end{lemma}

\begin{lemma}[Updating an entry of the drop-value vector]
\lean{Complexity.dropValues_update}
\label{Complexity.dropValues_update}
\leanok
\uses{Complexity.dropValues}
\difficulty{3}
Let $n,p,z,s,u,f,q,t,a$ and $b$ be natural numbers, and let $r$ be an index from $0$ to
$8$. Replacing entry $r$ of $(p,n,z,s,u,f,q,t,a)$ by $b$ gives the vector whose entries
are respectively $b$ when their indices equal $r$, and otherwise are the original
entries.
\end{lemma}

\begin{lemma}[Bound on the drop-value vector]
\lean{Complexity.dropValues_bound}
\label{Complexity.dropValues_bound}
\leanok
\uses{Complexity.dropValues}
\difficulty{1}
Let $n,p,z,s,u,f,q,t,a,B$ be natural numbers, each of $n,p,z,s,u,f,q,t,a$ at most $B$.
Then every entry of the length-nine vector $(p,n,z,s,u,f,q,t,a)$ is at most $B$.
\end{lemma}

\begin{lemma}[Tail decoding run]
\lean{Complexity.drop_tail_decode}
\label{Complexity.drop_tail_decode}
\leanok
\uses{Complexity.DropHalt, Complexity.DropReach, Complexity.LogProg.AReach.halt, Complexity.LogProg.AReach.refl, Complexity.LogProg.AReach.step, Complexity.LogProg.AReach.trans, Complexity.LogProg.Canon, Complexity.LogProg.astep, Complexity.LogProg.bitsVal, Complexity.LogProg.boundedDecode_loop, Complexity.LogProg.canonicalTail_loop, Complexity.LogProg.drop_canon, Complexity.LogProg.register_update_le, Complexity.dropARM, Complexity.dropCopy, Complexity.dropDecodeAdvance, Complexity.dropOracle, Complexity.dropTailAdvance, Complexity.dropValues, Complexity.dropValues_apply, Complexity.dropValues_bound, Complexity.drop_bit, Complexity.drop_copy_run, Complexity.drop_preS, Complexity.drop_ret, Complexity.pairEncode_getD, Complexity.query_accumulator_bound, Turing.length_pairEncode, Turing.pairEncode}
\difficulty{7}
Let $z$ and $t$ be finite bit strings, let $h=\langle z,t\rangle$, and set $p=2|z|+2$.
For either Boolean mode, if $t$ is canonical and $|t|$ is at most the length of the
binary representation of $|h|$, the machine has a run from the configuration with state
$\mathrm{copy}(0,\mathrm{false})$, pointer $p$, all other registers zero, and no answer
set, to the configuration with state $\mathrm{copy}(3,\mathrm{false})$, the same
pointer, query register $\mathrm{bitsVal}(t)$, tail register $p$, all other registers
zero, and no answer set. Every configuration before the endpoint satisfies the machine's
safety condition and has each register at most $4|h|+4$. If either condition on $t$
fails, the machine instead halts with answer false from that initial configuration, with
the same safety condition and register bound at every configuration before halting.
\end{lemma}

\begin{lemma}[Seek procedure run]
\lean{Complexity.drop_seek_run}
\label{Complexity.drop_seek_run}
\leanok
\uses{Complexity.DropReach, Complexity.LogProg.indexedSeek_loop, Complexity.dropARM, Complexity.dropOracle, Complexity.dropWalk}
\difficulty{4}
Let $h$ be a bit string, let $v$ assign a natural number to each of nine registers, and
let $\mathrm{mode}$ and $\mathrm{suffix}$ be Boolean parameters. Suppose every register
value is at most $4|h|+4$, register $2$ (if $\mathrm{suffix}$ is true) or register $5$
(otherwise) equals register $0$ plus $(2$ if true, $4$ otherwise$)\,\mathrm{count}$,
register $7+\mathrm{count}\le 4|h|+4$, and register $7\le$ register $6$. If register
$6<$ register $7+\mathrm{count}$, the machine reaches the hit state with register $0$
increased by $(2$ if $\mathrm{suffix}$ is true, $4$ otherwise$)(\text{register
}6-\text{register }7)$ and register $7$ set to register $6$. Otherwise it reaches the
answer-false state when $\mathrm{suffix}$ is true, or the clear-pointer stage $4$ when
it is false; in either case register $0$ becomes the initial value of register $2$ or
$5$, respectively, and register $7$ becomes its initial value plus $\mathrm{count}$. In
both cases the result register remains unset, and throughout the run every intermediate
configuration before the final one satisfies the machine's safety condition and has all
register values at most $4|h|+4$.
\end{lemma}

\begin{lemma}[Lookup from the Search State]
\lean{Complexity.drop_lookup}
\label{Complexity.drop_lookup}
\leanok
\uses{Complexity.DropHalt, Complexity.DropReach, Complexity.LogProg.AReach.halt, Complexity.LogProg.AReach.step, Complexity.LogProg.AReach.trans, Complexity.LogProg.astep, Complexity.LogProg.register_update_le, Complexity.dropARM, Complexity.dropCopy, Complexity.dropIdx, Complexity.dropIdx_pairEncode, Complexity.dropIdx_prefix_address, Complexity.dropIdx_suffix_address, Complexity.drop_copy_run, Complexity.drop_hit, Complexity.drop_preS, Complexity.drop_ret, Complexity.drop_seek_run, Turing.length_pairEncode, Turing.pairEncode}
\difficulty{7}
Let $mode$ be a Boolean mode, let $w,b,b'$ be bit strings, let $q$ be a natural number,
and let $v$ assign natural numbers to the nine registers. Suppose every register value
is at most $4|\langle\langle\langle w,b\rangle,b'\rangle,\mathrm{bits}(q)\rangle|+4$,
and the registers satisfy $v(0)=0$, $v(5)=4(2|w|+2)$, $v(3)=4|\langle w,b\rangle|+4$,
$v(2)=4|\langle w,b\rangle|+4+2|b'|$, $v(6)=q$, and $v(7)=0$. Starting on input
$\langle\langle\langle w,b\rangle,b'\rangle,\mathrm{bits}(q)\rangle$ in state
$seek(false)$ with values $v$ and no answer, the machine halts with the test
$q<|dropIdx(\langle\langle w,b\rangle,b'\rangle)|$ when $mode$ is true, and otherwise
with the bit of that decoded string at position $q$, defaulting to false when out of
range. Every configuration before halting satisfies the machine's safety condition and
register bound.
\end{lemma}

\begin{lemma}[First Component Read from a Pair]
\lean{Complexity.drop_first_read}
\label{Complexity.drop_first_read}
\leanok
\uses{Complexity.pairEncode_getD, Turing.pairEncode}
\difficulty{3}
Let $a$ and $b$ be bit strings, and let $i<|a|$. In the self-delimiting pairing of $a$
with $b$, the bit at position $2i$ is the bit of $a$ at position $i$ (with positions
counted from zero).
\end{lemma}

\begin{lemma}[Repeated Register Decrement Run]
\lean{Complexity.drop_decrement_run}
\label{Complexity.drop_decrement_run}
\leanok
\uses{Complexity.DropReach, Complexity.DropState, Complexity.LogProg.AReach.refl, Complexity.LogProg.AReach.step, Complexity.LogProg.astep, Complexity.LogProg.register_update_le, Complexity.dropARM, Complexity.drop_preS}
\difficulty{4}
Let $h$ be a bit string, let $r$ be one of nine registers, let $count$ be a natural
number, let $path : \mathbb{N} \to \mathrm{DropState}$ be a sequence of machine states,
and let $v$ assign a natural number to each register. Suppose every register initially
holds at most $4|h|+4$, where $|h|$ is the length of $h$, and at each step $i<count$ the
instruction in state $path(i)$ decrements register $r$ and moves to $path(i+1)$. Then
the machine in mode $mode$, started with state $path(0)$, register values $v$, and no
answer, reaches state $path(count)$ with register $r$ decreased by truncated subtraction
of $count$, all other registers unchanged, and no answer; throughout the run, every
preterminal configuration satisfies the machine's safety condition and every register
remains at most $4|h|+4$.
\end{lemma}

\begin{lemma}[Prepare the outer word]
\lean{Complexity.drop_z_prepare}
\label{Complexity.drop_z_prepare}
\leanok
\uses{Complexity.DropReach, Complexity.LogProg.AReach.refl, Complexity.LogProg.AReach.step, Complexity.LogProg.AReach.trans, Complexity.LogProg.astep, Complexity.dropARM, Complexity.dropCopy, Complexity.dropValues, Complexity.dropValues_bound, Complexity.drop_copy_run, Complexity.drop_decrement_run, Complexity.drop_preS}
\difficulty{5}
Let $\mathrm{mode}$ be Boolean, let $h$ be a bit string, and let $zlen$ and $q$ be
natural numbers satisfying $2zlen+2\le |h|$ and $q\le 2|h|$. Starting with the machine
in copy stage $3$ with pointer and tail both $2zlen+2$, query value $q$, and all other
registers zero except the second register, which is $|h|$, it reaches scan stage $1$
with pointer zero, second register $|h|$, end-marker register $2zlen$, query value $q$,
and tail value $2zlen+2$; all remaining registers are zero. The result remains unset,
and every intermediate configuration before the final one satisfies the machine's safety
condition and has all register values at most $4|h|+4$.
\end{lemma}

\begin{lemma}[Prepare the inner word]
\lean{Complexity.drop_u_prepare}
\label{Complexity.drop_u_prepare}
\leanok
\uses{Complexity.DropReach, Complexity.LogProg.AReach.refl, Complexity.LogProg.AReach.step, Complexity.LogProg.AReach.trans, Complexity.LogProg.astep, Complexity.dropARM, Complexity.dropCopy, Complexity.dropValues, Complexity.dropValues_bound, Complexity.drop_copy_run, Complexity.drop_decrement_run, Complexity.drop_preS}
\difficulty{5}
Let $\mathrm{mode}$ be Boolean, let $h$ be a bit string, and let $zlen$, $ulen$, and $q$
be natural numbers satisfying $2zlen+2\le |h|$, $2ulen+2\le zlen$, and $q\le 2|h|$.
Starting with the machine in copy stage $4$ with pointer $4ulen+4$, the second register
equal to $|h|$, query register $q$, tail register $2zlen+2$, and all other registers
zero, the machine reaches scan stage $2$ with those pointer, query, and tail values
changed to zero, $q$, and $2zlen+2$, respectively; the second register remains $|h|$,
the retained register becomes $4ulen+4$, and the inner register becomes $4ulen$. The
result remains unset, and every intermediate configuration before the final one
satisfies the machine's safety condition and has all register values at most $4|h|+4$.
\end{lemma}

\begin{lemma}[Prepare Registers for Lookup]
\lean{Complexity.drop_lookup_prepare}
\label{Complexity.drop_lookup_prepare}
\leanok
\uses{Complexity.DropReach, Complexity.LogProg.AReach.refl, Complexity.LogProg.AReach.step, Complexity.LogProg.AReach.trans, Complexity.LogProg.astep, Complexity.dropARM, Complexity.dropCopy, Complexity.dropValues, Complexity.dropValues_bound, Complexity.drop_copy_run, Complexity.drop_preS}
\difficulty{4}
Let $mode$ be a Boolean mode, let $h$ be a bit string, and let $zlen, ulen, wlen, q$ be
natural numbers such that $2zlen+2\le |h|$, $2ulen+2\le zlen$, $2wlen+2\le ulen$, and
$q\le 2|h|$. Starting with state $copy(6,false)$, no answer, and register values
$ptr=8wlen+8$, $zend=2zlen$, $retained=4ulen+4$, $inner=4ulen$, $query=q$,
$tail=2zlen+2$, and all remaining registers zero, the machine reaches state
$seek(false)$ with no answer and register values $zend=2zlen$, $retained=4ulen+4$,
$inner=4ulen$, $kept=8wlen+8$, $query=q$, $tail=2zlen+2$, and all remaining registers
zero. Every configuration before the endpoint satisfies the machine's safety condition
and register bound.
\end{lemma}

\begin{lemma}[Inner Lookup Run]
\lean{Complexity.drop_inner_run}
\label{Complexity.drop_inner_run}
\leanok
\uses{Complexity.DropHalt, Complexity.DropReach, Complexity.LogProg.AReach.halt, Complexity.LogProg.AReach.trans, Complexity.LogProg.register_update_le, Complexity.dropIdx, Complexity.dropLimit, Complexity.dropScanDone, Complexity.dropStride, Complexity.dropValues, Complexity.dropValues_apply, Complexity.dropValues_bound, Complexity.drop_first_read, Complexity.drop_lookup, Complexity.drop_lookup_prepare, Complexity.drop_pair_run, Complexity.drop_ret, Complexity.drop_u_prepare, Turing.eq_pairEncode_of_pairDecode, Turing.length_pairEncode, Turing.pairDecode, Turing.pairDecode_pairEncode, Turing.pairEncode}
\difficulty{6}
Let $mode$ be a Boolean mode, let $u$ and $b'$ be bit strings, and let $q$ be a natural
number satisfying $q\le 2|\langle\langle u,b'\rangle,\mathrm{bits}(q)\rangle|$, where
$\mathrm{bits}(q)$ is the binary bit string for $q$. Set $h=\langle\langle
u,b'\rangle,\mathrm{bits}(q)\rangle$ and $z=\langle u,b'\rangle$. Starting on input $h$
in state $copy(4,false)$, with register values $ptr=4|u|+4$, $zend=2|z|$, $query=q$,
$tail=2|z|+2$, and all other registers zero, the machine halts with the length test
$q<|dropIdx(z)|$ when $mode$ is true, and otherwise with the bit of $dropIdx(z)$ at
position $q$, defaulting to false if out of range. Every configuration before halting
satisfies the machine's safety condition and register bound.
\end{lemma}

\begin{lemma}[Nested Drop Computation]
\lean{Complexity.drop_nested_run}
\label{Complexity.drop_nested_run}
\leanok
\uses{Complexity.DropHalt, Complexity.DropReach, Complexity.LogProg.AReach.halt, Complexity.LogProg.AReach.trans, Complexity.LogProg.register_update_le, Complexity.dropIdx, Complexity.dropLimit, Complexity.dropScanDone, Complexity.dropStride, Complexity.dropValues, Complexity.dropValues_bound, Complexity.drop_first_read, Complexity.drop_inner_run, Complexity.drop_pair_run, Complexity.drop_ret, Complexity.drop_z_prepare, Turing.eq_pairEncode_of_pairDecode, Turing.length_pairEncode, Turing.pairDecode, Turing.pairDecode_pairEncode, Turing.pairEncode}
\difficulty{6}
Let $m\in\{\mathrm{false},\mathrm{true}\}$, let $z$ be a bit string, and let
$q\in\mathbb N$ satisfy $q\le 2|\langle z,\mathrm{bits}(q)\rangle|$. On input $h=\langle
z,\mathrm{bits}(q)\rangle$, the specified computation starting in its copy state, with
pointer register $2|z|+2$, query register $q$, tail register $2|z|+2$, all other
registers zero, and no answer, halts with the answer that is true exactly when either
$q$ is less than the length of the transformed word (if $m$ is true) or its $q$-th bit
is $1$ (if $m$ is false). Every configuration before halting satisfies the stated
computation invariant and has each register at most $4|h|+4$.
\end{lemma}

\begin{lemma}[Language Indicator on Encoded Pairs]
\lean{Complexity.drop_lang_pair}
\label{Complexity.drop_lang_pair}
\leanok
\uses{Complexity.dropIdx, Complexity.dropLang, Complexity.indicator_indexLang_pairEncode, Turing.MultiTapeTM.indicator, Turing.pairEncode}
\difficulty{1}
For every Boolean mode, bit string $z$, and natural number $q$, the Boolean indicator of
the language that tests whether $q$ is less than the length of $dropIdx(z)$ in length
mode, and otherwise tests whether the bit of $dropIdx(z)$ at position $q$ is true,
evaluated on $\langle z,\mathrm{bits}(q)\rangle$, equals true or false according to that
same test.
\end{lemma}

\begin{lemma}[Malformed input is rejected]
\lean{Complexity.drop_bad_shape}
\label{Complexity.drop_bad_shape}
\leanok
\uses{Complexity.dropLang, Turing.MultiTapeTM.indicator, Turing.pairDecode, Turing.pairDecode_pairEncode}
\difficulty{4}
For either choice of the mode, let $h$ be a bit string that the aligned pair parser
cannot decode. The Boolean indicator of the language that tests the transformed input is
false on $h$.
\end{lemma}

\begin{lemma}[Noncanonical tail is rejected]
\lean{Complexity.drop_bad_tail}
\label{Complexity.drop_bad_tail}
\leanok
\uses{Complexity.LogProg.Canon, Complexity.LogProg.canon_bits, Complexity.dropLang, Turing.MultiTapeTM.indicator, Turing.pairEncode, Turing.pairEncode_injective}
\difficulty{4}
For either choice of the mode, let $z$ and $tail$ be bit strings, and suppose the last
bit of every nonempty canonical bit string condition fails for $tail$ (that is, $tail$
is not canonical). The Boolean indicator of the transformed-input language is false on
the encoding of $z$ and $tail$.
\end{lemma}

\begin{lemma}[Rejection of Invalid Trailing Payloads]
\lean{Complexity.drop_rejected_tail}
\label{Complexity.drop_rejected_tail}
\leanok
\uses{Complexity.LogProg.Canon, Complexity.LogProg.bitsVal, Complexity.LogProg.canon_eq_bits, Complexity.dropIdx_query_overflow, Complexity.dropLang, Complexity.drop_bad_tail, Complexity.drop_lang_pair, Turing.MultiTapeTM.indicator, Turing.pairEncode}
\difficulty{5}
Let $m\in\{\mathrm{false},\mathrm{true}\}$ and let $z$ and $t$ be bit strings. If $t$
either is empty or does not end in $1$, or if $|t|$ exceeds the length of the binary
representation of the length of $\langle z,t\rangle$, then the indicator of the mode-$m$
drop language is false on $\langle z,t\rangle$.
\end{lemma}

\begin{lemma}[Drop Computation Correctness]
\lean{Complexity.drop_run}
\label{Complexity.drop_run}
\leanok
\uses{Complexity.DropHalt, Complexity.DropReach, Complexity.LogProg.AReach.halt, Complexity.LogProg.AReach.trans, Complexity.LogProg.Canon, Complexity.LogProg.bitsVal, Complexity.LogProg.bitsVal_bits, Complexity.LogProg.canon_eq_bits, Complexity.LogProg.inputLength_loop, Complexity.LogProg.register_update_le, Complexity.dropARM, Complexity.dropLang, Complexity.dropLimit, Complexity.dropOracle, Complexity.dropScanDone, Complexity.dropStride, Complexity.dropValues, Complexity.dropValues_bound, Complexity.drop_bad_shape, Complexity.drop_lang_pair, Complexity.drop_len, Complexity.drop_nested_run, Complexity.drop_pair_run, Complexity.drop_rejected_tail, Complexity.drop_ret, Complexity.drop_tail_decode, Complexity.query_index_bound, Turing.MultiTapeTM.indicator, Turing.eq_pairEncode_of_pairDecode, Turing.pairDecode, Turing.pairDecode_pairEncode, Turing.pairEncode}
\difficulty{6}
For every mode $m\in\{\mathrm{false},\mathrm{true}\}$ and every input bit string $h$,
the specified computation started in its length-test state with all nine registers zero
and no answer halts with the membership indicator of the mode-$m$ drop language on $h$.
Every configuration before halting satisfies the stated computation invariant.
\end{lemma}

\begin{lemma}[Logarithmic Space for the Drop Languages]
\lean{Complexity.drop_mem}
\label{Complexity.drop_mem}
\leanok
\uses{Complexity.LOGSPACE, Complexity.LogProg.AHalt.mono, Complexity.LogProg.arm_decides_poly, Complexity.dropARM, Complexity.dropBank, Complexity.dropLang, Complexity.drop_run, Complexity.inputBitLang_mem, Complexity.inputLenLang_mem}
\difficulty{4}
For each choice of mode $m\in\{\mathrm{false},\mathrm{true}\}$, the language defined by
accepting an encoded pair $(z,q)$ when either $q$ is less than the length of the
transformed word obtained by dropping the nested encoding from $z$ (if
$m=\mathrm{true}$), or its $q$-th bit is $1$ (if $m=\mathrm{false}$), is decidable in
logarithmic space.
\end{lemma}

\begin{theorem}[Logspace computability of index deletion]
\lean{Complexity.implicitlyLogspaceComputable_dropIdx}
\label{Complexity.implicitlyLogspaceComputable_dropIdx}
\leanok
\uses{Complexity.ImplicitlyLogspaceComputable, Complexity.dropIdx, Complexity.dropIdx_length_le, Complexity.dropLang, Complexity.drop_mem}
\difficulty{3}
For a bit string $z$, parse it as an encoded pair $(u,b')$; then parse $u$ as an encoded
pair $(w,b'')$ and return the encoding of $(w,b')$. If either parse fails, return the
empty string. This resulting function is implicitly logspace computable: its output
length is bounded by $C(|x|+1)^c$ for some natural numbers $C,c$ on every input $x$, and
both the language of encoded pairs $(x,\text{bits}(i))$ for which output bit $i$ is $1$
and the language of such pairs for which $i$ is less than the output length are
decidable in logarithmic space.
\end{theorem}

\begin{definition}[Languages encoding a computation]
\lean{Complexity.compLift}
\label{Complexity.compLift}
\leanok
\uses{Complexity.dropIdx, Complexity.indexLang}
Let $f$ be a function from finite Boolean lists to finite Boolean lists, and let $m$ be
a Boolean flag. Define a language of Boolean lists consisting of those $z$ for which
deleting the first two entries of $z$ yields a list in the indexed language whose
predicate at input $x$ and index $i$ holds exactly when either $m$ is true and
$i<|f(x)|$, or $m$ is false and the $i$th entry of $f(x)$ (taken to be false when the
index is out of range) is true.
\end{definition}

\begin{lemma}[Logspace membership of a lifted query language]
\lean{Complexity.compLift_mem}
\label{Complexity.compLift_mem}
\leanok
\uses{Complexity.ImplicitlyLogspaceComputable, Complexity.LOGSPACE, Complexity.compLift, Complexity.implicitlyLogspaceComputable_dropIdx, Complexity.indexLang, Complexity.mem_LOGSPACE_preimage}
\difficulty{3}
Let $f$ map finite Boolean strings to finite Boolean strings and suppose its output
length is polynomially bounded in the input length, with both the indexed output-bit
language and the indexed output-length language decidable in logarithmic space. For
either choice of query mode, the language obtained by applying the two-stage
index-dropping transformation and then testing the corresponding output bit or output
length is decidable in logarithmic space.
\end{lemma}

\begin{definition}[Four-language computation bank]
\lean{Complexity.compBank}
\label{Complexity.compBank}
\leanok
\uses{Complexity.compLift, Complexity.inputBitLang, Complexity.inputLenLang}
Let $f$ be a function from finite Boolean lists to finite Boolean lists. Define a
function on the four-element set $\{0,1,2,3\}$ with values in languages of Boolean lists
by assigning, in order, the languages for input bits, input length, the computation
language for $f$ with length mode false, and the computation language for $f$ with
length mode true.
\end{definition}

\begin{lemma}[Logspace membership of the query bank]
\lean{Complexity.compBank_mem}
\label{Complexity.compBank_mem}
\leanok
\uses{Complexity.ImplicitlyLogspaceComputable, Complexity.LOGSPACE, Complexity.compBank, Complexity.compLift_mem, Complexity.inputBitLang_mem, Complexity.inputLenLang_mem}
\difficulty{4}
Let $f$ map finite Boolean strings to finite Boolean strings. Suppose its output length
is bounded by a polynomial in the input length, and that both indexed languages—testing
whether the $i$th output bit is true and testing whether the output has length greater
than $i$—are decidable in logarithmic space. Then, for each of the four languages in the
bank consisting of the input-bit query, the input-length query, and the corresponding
two queries for $f$, that language is decidable in logarithmic space.
\end{lemma}

\begin{lemma}[Value of a lifted query]
\lean{Complexity.compLift_query}
\label{Complexity.compLift_query}
\leanok
\uses{Complexity.compLift, Complexity.dropIdx_pairEncode, Complexity.indexLang, Complexity.indicator_indexLang_pairEncode, Turing.MultiTapeTM.indicator, Turing.pairEncode}
\difficulty{4}
Let $f$ map finite Boolean strings to finite Boolean strings, let $x,b$ be finite
Boolean strings, and let $p,j$ be natural numbers. On the encoding formed by pairing $x$
with $b$, then pairing the result with the binary representation of $p$, and then
pairing that result with the binary representation of $j$, the Boolean indicator of the
lifted query language equals the following: in length mode, it is true exactly when
$j<|f(x)|$; otherwise, it is the $j$th bit of $f(x)$, with false used when that position
is beyond the end of the string.
\end{lemma}

\begin{lemma}[Output length bound for the query transformation]
\lean{Complexity.compQ_host_cap}
\label{Complexity.compQ_host_cap}
\leanok
\uses{Complexity.compQ, Complexity.compQ_length_le, Turing.length_pairEncode, Turing.pairEncode}
\difficulty{3}
Let $f$ map finite Boolean strings to finite Boolean strings, and suppose there are
natural numbers $C,e$ such that $|f(x)|\le C(|x|+1)^e$ for every finite Boolean string
$x$. For any finite Boolean strings $w$ and $t$, the output length of the query
transformation applied to $w$ is at most $(2C+3)(|\langle w,t\rangle|+1)^{e+1}$, where
$\langle w,t\rangle$ denotes their self-delimiting pair encoding.
\end{lemma}

\begin{lemma}[Overflowing query is out of range]
\lean{Complexity.compQ_query_overflow}
\label{Complexity.compQ_query_overflow}
\leanok
\uses{Complexity.LogProg.length_bits_mono, Complexity.compQ, Complexity.compQ_host_cap, Turing.pairEncode}
\difficulty{3}
Let $f$ map bit strings to bit strings, let $C,e$ be natural numbers, and suppose
$|f(x)|\le C(|x|+1)^e$ for every bit string $x$. For bit strings $w$ and a natural
number $p$, suppose the binary expansion of $(2C+3)(|\langle
w,\operatorname{bin}(p)\rangle|+1)^{e+1}$ is shorter than the binary expansion of $p$.
Then the output bit at position $p$ of the transformed query on $w$ is false, and $p$ is
not less than the length of that output.
\end{lemma}

\begin{definition}[Computation machine states]
\lean{Complexity.CompState}
\label{Complexity.CompState}
\leanok
For natural numbers $C$ and $e$, define an inductive type of states with the following
constructors: the parameter-free states listed as count, count increment, base
increment, tail check, tail read, tail done, reset decode, decode, decode down, decode
read, decode done, canonical, canonical reset, canonical read, inner prepare, length
bound, length query, length increment, double test, double down, lookup reset, prefix
hit, prefix bit, suffix hit, and suffix reset; states additionally parameterized by
elements of the indicated finite sets and Boolean values as specified in their
constructors, including copy stages in an eight-element set, counters in an
$(e+2)$-element set, and the factor-increment counter in a $(2C+3)$-element set.
\end{definition}

\begin{definition}[Copy-stage transition table]
\lean{Complexity.compCopy}
\label{Complexity.compCopy}
\leanok
\uses{Complexity.CompState}
For natural numbers $C$ and $e$, define a function on the eight stages $0,1,2,3,4,5,6,7$
whose values are triples of elements of $\text{Fin}(14)$ and computation states, in
order: $(2,0,\text{the copy state at stage }1\text{ with bump false})$, $(3,2,\text{the
outer-decrement state with second flag false})$, $(8,1,\text{the base-increment state})$,
$(0,1,\text{the decode state})$, $(4,0,\text{the canonical state})$, $(0,3,\text{the
canonical-decrement state with second flag false})$, $(8,4,\text{the inner-halving state
with second flag false})$, and $(0,4,\text{the seek state with suffix flag true})$.
\end{definition}

\begin{definition}[Stride determined by an inner flag]
\lean{Complexity.compStride}
\label{Complexity.compStride}
\leanok
For a Boolean parameter $b$, define the stride to be $2$ when $b$ is true and $1$ when
$b$ is false.
\end{definition}

\begin{definition}[Limit determined by an inner flag]
\lean{Complexity.compLimit}
\label{Complexity.compLimit}
\leanok
For a Boolean parameter $b$, define an element of the fourteen-element set to be $3$
when $b$ is true and $1$ when $b$ is false.
\end{definition}

\begin{definition}[Advancing a computation state]
\lean{Complexity.compAdvance}
\label{Complexity.compAdvance}
\leanok
\uses{Complexity.CompState, Complexity.compStride}
For natural numbers $C,e$, a Boolean parameter $b$, an element $c$ of the four-element
set, and a natural number $i$, define a computation state as follows. If $i$ is less
than the stride determined by $b$, the state is the scan-increment state with parameters
$b,c,i$; otherwise it is selected from the four-entry list consisting of the second
state with false, the second state with true, the scan state, and the copy state at
stage $4$ when $b$ is true and stage $0$ when $b$ is false, with false bump flag, using
$c$ as the index.
\end{definition}

\begin{definition}[Complexity factor]
\lean{Complexity.compFactor}
\label{Complexity.compFactor}
\leanok
For a natural number $C$ and a Boolean flag $o$, define the factor to be $2C+3$ when $o$
is true and $C$ when $o$ is false.
\end{definition}

\begin{definition}[Factor advance state]
\lean{Complexity.compFactorAdvance}
\label{Complexity.compFactorAdvance}
\leanok
\uses{Complexity.CompState, Complexity.compFactor}
Given natural numbers $C,e,i$ and a Boolean flag $o$, define the state to be the
factor-increment state with counter $i$ when $i$ is less than the complexity factor for
$C,o$. Otherwise, define it to be the factor-test state for $o$.
\end{definition}

\begin{definition}[Decode advance state]
\lean{Complexity.compDecodeAdvance}
\label{Complexity.compDecodeAdvance}
\leanok
\uses{Complexity.CompState, Complexity.dropDecodeAdvance}
Given natural numbers $C,e,i$ and a Boolean $b$, apply the decode-advance operation for
$b$ to $i$. If its result is a decode-increment state, return that state; in every other
case, return the decode-scale state for $b$.
\end{definition}

\begin{definition}[Complexity machine instruction map]
\lean{Complexity.compARM}
\label{Complexity.compARM}
\leanok
\uses{Complexity.CompState, Complexity.LogProg.ARM, Complexity.compAdvance, Complexity.compCopy, Complexity.compFactorAdvance, Complexity.compLimit}
For natural numbers $C,e$ and a Boolean mode, define a deterministic instruction map on
the states of a finite-register machine. The machine has $14$ registers and $4$ oracle
deciders; its instructions are determined by the state, the parameters, and the mode as
specified by the transition cases.
\end{definition}

\begin{definition}[Oracle-language indicator]
\lean{Complexity.compOracle}
\label{Complexity.compOracle}
\leanok
\uses{Complexity.compBank, Turing.MultiTapeTM.indicator}
Given a function $f$ from Boolean lists to Boolean lists, an index $j$ among four oracle
deciders, and a Boolean list $w$, define the oracle value to be the indicator of whether
$w$ belongs to the language assigned to $f$ and $j$ by the oracle bank construction.
\end{definition}

\begin{definition}[Complexity bound]
\lean{Complexity.compBound}
\label{Complexity.compBound}
\leanok
For natural numbers $C$, $e$, and $n$, define the bound $B(C,e,n)=8(2C+3)(n+1)^{e+1}$.
\end{definition}

\begin{lemma}[Machine instruction admissibility]
\lean{Complexity.comp_preS}
\label{Complexity.comp_preS}
\leanok
\uses{Complexity.CompState, Complexity.LogProg.AConf, Complexity.LogProg.PreS, Complexity.compARM, Complexity.compCopy, Complexity.compLimit}
\difficulty{5}
For all natural numbers $C,e$, every Boolean value $m$, every finite bit string $h$, and
every configuration $a$ consisting of an optional state, fourteen natural-number
register values, and an optional Boolean result, the instruction specified at its active
state satisfies its admissibility condition: a register-equality instruction compares
distinct registers; an input-word comparison is used only when the input has the
required plain or nested canonical encoding; and an oracle call uses a list of pairwise
distinct argument registers.
\end{lemma}

\begin{lemma}[Input-bit oracle answer]
\lean{Complexity.comp_bit}
\label{Complexity.comp_bit}
\leanok
\uses{Complexity.compOracle, Complexity.drop_bit, Turing.pairEncode}
\difficulty{1}
Let $f$ map bit strings to bit strings, let $h$ be a bit string, and let $i$ be a
natural number. The first oracle in the four-language bank associated with $f$, queried
on the encoding of $(h,\operatorname{bin}(i))$, returns the bit of $h$ at position $i$,
taking the value false when that position is beyond the end of $h$.
\end{lemma}

\begin{lemma}[Input-length oracle]
\lean{Complexity.comp_len}
\label{Complexity.comp_len}
\leanok
\uses{Complexity.compOracle, Complexity.drop_len, Turing.pairEncode}
\difficulty{1}
For every bit string $h$ and natural number $i$, the length oracle on the encoded pair
of $h$ and the binary representation of $i$ returns true exactly when $i<|h|$, where
$|h|$ is the length of $h$.
\end{lemma}

\begin{lemma}[Lifted output oracle answers]
\lean{Complexity.comp_lifted}
\label{Complexity.comp_lifted}
\leanok
\uses{Complexity.compLift_query, Complexity.compOracle, Turing.pairEncode}
\difficulty{4}
Let $f$ be a function on bit strings, let $m$ be a Boolean mode, let $x$ and $b$ be bit
strings, and let $p,i$ be natural numbers. The lifted oracle, queried at the encoding
formed by nesting the pairs $(x,b)$, then its binary index $p$, then the binary index
$i$, returns true exactly when $i<|f(x)|$ if $m$ is true; if $m$ is false, it returns
the bit of $f(x)$ at position $i$, with out-of-range positions interpreted as false.
\end{lemma}

\begin{definition}[Good machine configuration]
\lean{Complexity.CompGood}
\label{Complexity.CompGood}
\leanok
\uses{Complexity.CompState, Complexity.LogProg.AConf, Complexity.LogProg.PreS, Complexity.compARM}
For natural numbers $C,e,B$, a Boolean mode, a Boolean-list input $h$, and a
configuration $a$ of the machine with $14$ registers, call $a$ good when it is a
pre-configuration for the specified machine on input $h$ and every register value in $a$
is at most $B$.
\end{definition}

\begin{definition}[Bounded computational reachability]
\lean{Complexity.CompReach}
\label{Complexity.CompReach}
\leanok
\uses{Complexity.CompGood, Complexity.LogProg.AReach, Complexity.compARM, Complexity.compOracle}
For natural numbers $C,e,B$, a Boolean mode, a function $f$ from Boolean lists to
Boolean lists, and an input $h$ that is a Boolean list, define the bounded reachability
predicate to mean reachability from $h$ for the machine determined by $C,e,$ and the
mode, using the oracle determined by $f$, subject to the good-configuration condition
determined by $C,e,$ the mode, $h$, and $B$.
\end{definition}

\begin{definition}[Complexity-bounded oracle halting]
\lean{Complexity.CompHalt}
\label{Complexity.CompHalt}
\leanok
\uses{Complexity.CompGood, Complexity.LogProg.AHalt, Complexity.compARM, Complexity.compOracle}
For natural numbers $C,e,B$, a Boolean mode $m$, a function $f$ from Boolean lists to
Boolean lists, and an initial Boolean list $h$, write $P(C,e,m,f,h,B)$ for the assertion
that the computation determined by $C,e,m$, with oracle $f$ and initial input $h$, halts
and satisfies the associated goodness condition with bound $B$.
\end{definition}

\begin{lemma}[Product Loop Register Update]
\lean{Complexity.comp_product}
\label{Complexity.comp_product}
\leanok
\uses{Complexity.CompReach, Complexity.LogProg.AReach.single, Complexity.LogProg.AReach.step, Complexity.LogProg.AReach.trans, Complexity.LogProg.astep, Complexity.LogProg.registerCopy_loop, Complexity.LogProg.registerScale_loop, Complexity.LogProg.register_update_le, Complexity.compARM, Complexity.compOracle, Complexity.comp_preS}
\difficulty{6}
Let $C,e,B\in\mathbb{N}$, let $\mathrm{mode}$ and $\mathrm{outer}$ be Boolean values,
let $f$ map finite Boolean strings to finite Boolean strings, let $h$ be a finite
Boolean string, let $k\in\mathrm{Fin}(e+2)$, and let $v$ be a vector of fourteen
natural-number registers. Suppose every register of $v$ is at most $B$, register $6$ is
zero, $v(10)\le v(8)$, and $v(7)+(v(8)-v(10))v(11)\le B$. The computation can reach a
configuration whose state is the product-clean state with its second flag false, whose
register vector is obtained from $v$ by setting register $10$ to $v(8)$ and register $7$
to $v(7)+(v(8)-v(10))v(11)$, and whose result is still absent. Every configuration
strictly before this endpoint satisfies the machine precondition and has all registers
at most $B$.
\end{lemma}

\begin{lemma}[Register multiplication run]
\lean{Complexity.comp_multiply}
\label{Complexity.comp_multiply}
\leanok
\uses{Complexity.CompReach, Complexity.LogProg.AReach.single, Complexity.LogProg.AReach.step, Complexity.LogProg.AReach.trans, Complexity.LogProg.PreS, Complexity.LogProg.astep, Complexity.LogProg.registerCopy_loop, Complexity.LogProg.register_update_le, Complexity.compARM, Complexity.compOracle, Complexity.comp_preS, Complexity.comp_product}
\difficulty{6}
Let $C,e$ be natural numbers, let $m$ and $o$ be Booleans, let $f$ map finite bit
strings to finite bit strings, let $h$ be a finite bit string, let $k$ be an element of
$\text{Fin}(e+2)$, let $B$ be a natural number, and let $v$ assign a natural number to
each of fourteen registers. Suppose every register value is at most $B$, registers $6$,
$10$, and $11$ are $0$, and the product of the values in registers $7$ and $8$ is at
most $B$. Then the computation can run from the product-copy state with parameters $o,k$
and register vector $v$ to the power state with parameters $o,k-1$ and register vector
obtained by replacing register $7$ with the product of its initial value and register
$8$, while the result remains empty and every configuration before the endpoint
satisfies the program’s instruction-side conditions and register bound $B$.
\end{lemma}

\begin{lemma}[Factor multiplication run]
\lean{Complexity.comp_factor_run}
\label{Complexity.comp_factor_run}
\leanok
\uses{Complexity.CompReach, Complexity.LogProg.AReach.step, Complexity.LogProg.AReach.trans, Complexity.LogProg.registerCopy_loop, Complexity.LogProg.registerScale_loop, Complexity.LogProg.register_update_le, Complexity.compARM, Complexity.compFactor, Complexity.compFactorAdvance, Complexity.compOracle, Complexity.comp_preS}
\difficulty{5}
Let $C,e,B$ be natural numbers, let $m$ and $o$ be Boolean flags, let $f$ map finite
Boolean words to finite Boolean words, and let $h$ be a finite Boolean word. Suppose $v$
assigns a natural number to each of fourteen registers, every register value is at most
$B$, register $11$ is zero, and $Fv_7\le B$, where $F=C$ if $o$ is false and $F=2C+3$ if
$o$ is true. Starting in the factor-copy state with flag $o$ and no result, the
computation reaches the length-bound state if $o$ is false and the tail-test state with
flag true if $o$ is true, with register $7$ equal to $Fv_7$ and every other register
unchanged. Every configuration before the final one satisfies the prescribed invariant:
the instruction is admissible at that configuration and all register values are at most
$B$.
\end{lemma}

\begin{lemma}[Register update in the power loop]
\lean{Complexity.comp_power_run}
\label{Complexity.comp_power_run}
\leanok
\uses{Complexity.CompReach, Complexity.LogProg.AReach.step, Complexity.LogProg.AReach.trans, Complexity.LogProg.astep, Complexity.LogProg.register_update_le, Complexity.compARM, Complexity.compFactor, Complexity.comp_factor_run, Complexity.comp_multiply, Complexity.comp_preS}
\difficulty{6}
Let $C,e,B\in\mathbb{N}$, let $m,o\in\{\mathrm{false},\mathrm{true}\}$, let $f$ map
finite bit strings to finite bit strings, let $h$ be a finite bit string, let $k$ be an
element of the finite set of natural numbers below $e+2$, and let $v$ be a vector of
fourteen natural-number registers. Suppose every register value is at most $B$,
registers $6,10,11$ are zero, register $8$ is positive, and, with $F=2C+3$ if $o$ is
true and $F=C$ otherwise, both $v(7)v(8)^{k}$ and $Fv(7)v(8)^{k}$ are at most $B$. Then
there is a run from the power-loop state with flags $o,k$, register vector $v$, and no
answer yet, to length bound if $o$ is false and tail test if $o$ is true. Every
configuration before the endpoint satisfies the program's admissibility condition and
has all register values at most $B$; at the endpoint register $7$ is $Fv(7)v(8)^k$, and
every other register retains its initial value.
\end{lemma}

\begin{lemma}[Polynomial counter run]
\lean{Complexity.comp_poly_run}
\label{Complexity.comp_poly_run}
\leanok
\uses{Complexity.CompReach, Complexity.LogProg.AReach.step, Complexity.LogProg.astep, Complexity.LogProg.register_update_le, Complexity.compARM, Complexity.compFactor, Complexity.comp_power_run, Complexity.comp_preS}
\difficulty{4}
Let $C,e,B\in\mathbb{N}$, let $m,o\in\{\mathrm{false},\mathrm{true}\}$, let $f$ map
finite bit strings to finite bit strings, let $h$ be a finite bit string, and let $v$ be
a vector of fourteen natural-number registers. Suppose every register value is at most
$B$, registers $6,7,10,11$ are zero, register $8$ is positive, and, with $q=e+1$ if $o$
is true and $q=e$ otherwise and $F=2C+3$ if $o$ is true and $F=C$ otherwise, both
$v(8)^q$ and $Fv(8)^q$ are at most $B$. Then there is a run of the specified register
program from state polynomial start with flag $o$ and register vector $v$ to tail test
if $o$ is true and to length bound otherwise. Every configuration before the endpoint
satisfies the program's admissibility condition and has all register values at most $B$;
at the endpoint, register $7$ is $Fv(8)^q$ and every other register retains its initial
value.
\end{lemma}

\begin{lemma}[Power and factor bounds]
\lean{Complexity.comp_power_bound}
\label{Complexity.comp_power_bound}
\leanok
\uses{Complexity.compBound, Complexity.compFactor}
\difficulty{3}
Let $C,e,n,b,r$ be natural numbers and let $o$ be a Boolean parameter. If $b\le n+1$ and
$r\le e+1$, then $b^r\le 8(2C+3)(n+1)^{e+1}$ and $F b^r\le 8(2C+3)(n+1)^{e+1}$, where
$F=2C+3$ if $o$ is true and $F=C$ otherwise.
\end{lemma}

\begin{lemma}[Room in the computational bound]
\lean{Complexity.comp_bound_room}
\label{Complexity.comp_bound_room}
\leanok
\uses{Complexity.compBound}
\difficulty{3}
For natural numbers $C,e,n$, let $B=8(2C+3)(n+1)^{e+1}$. Then $n+1\le B$ and
$2(2C+3)(n+1)^{e+1}\le B$.
\end{lemma}

\begin{lemma}[Length search register run]
\lean{Complexity.comp_length_search}
\label{Complexity.comp_length_search}
\leanok
\uses{Complexity.CompReach, Complexity.LogProg.AReach.refl, Complexity.LogProg.AReach.single, Complexity.LogProg.AReach.step, Complexity.LogProg.astep, Complexity.LogProg.register_update_le, Complexity.LogProg.vword_whole₁, Complexity.compARM, Complexity.compOracle, Complexity.comp_preS, Turing.pairEncode}
\difficulty{6}
Let $C,e,B,F$ be natural numbers, let $m$ be a Boolean, let $f$ map finite bit strings
to finite bit strings, let $h$ be a finite bit string, and let $v$ assign a natural
number to each of fourteen registers. Suppose every register value is at most $B$,
register $12$ has value at most $F$, and $F$ is at most the value in register $7$.
Assume that for every $i$ smaller than the value in register $7$, the length oracle for
$f$ on the encoded pair $(h,i)$ returns true exactly when $i<F$. Then the computation
can run from its length-bound state with register vector $v$ and no answer to its
double-test state with register $12$ set to $F$ and no answer, with every configuration
before the endpoint satisfying the program’s instruction-side conditions and register
bound $B$.
\end{lemma}

\begin{definition}[Complexity register values]
\lean{Complexity.compValues}
\label{Complexity.compValues}
\leanok
For natural numbers $n$, and natural numbers
$\mathit{ptr},\mathit{tail},\mathit{stop},\mathit{inner},\mathit{query},\mathit{temp},\mathit{cap},\mathit{base},\mathit{spare},\mathit{iter},\mathit{old},\mathit{len},\mathit{idx}$,
each defaulting to $0$, define a function on the indices $0,\ldots,13$ whose values in
order are those fourteen parameters, with $n$ in the second position.
\end{definition}

\begin{lemma}[Coordinate evaluation of the value vector]
\lean{Complexity.compValues_apply}
\label{Complexity.compValues_apply}
\leanok
\uses{Complexity.compValues}
\difficulty{1}
For natural numbers $n,p,a,w,b,q,t,c,x,s,i,y,l,j$ and an index $r\in\{0,1,\ldots,13\}$,
the $r$th coordinate of the value vector $(p,n,a,w,b,q,t,c,x,s,i,y,l,j)$ is the $r$th
coordinate of that displayed vector.
\end{lemma}

\begin{lemma}[Updating a coordinate of the value vector]
\lean{Complexity.compValues_update}
\label{Complexity.compValues_update}
\leanok
\uses{Complexity.compValues}
\difficulty{3}
For natural numbers $n,p,a,w,b,q,t,c,x,s,i,y,l,j,v$ and an index
$r\in\{0,1,\ldots,13\}$, updating the value vector $(p,n,a,w,b,q,t,c,x,s,i,y,l,j)$ at
coordinate $r$ to $v$ gives the vector whose entries, in order, are $p$ if $r\ne0$ and
$v$ otherwise; $n$ if $r\ne1$ and $v$ otherwise; and, for each $k=2,\ldots,13$, the
original entry at coordinate $k$ if $r\ne k$ and $v$ otherwise.
\end{lemma}

\begin{lemma}[Register copy run]
\lean{Complexity.comp_copy_run}
\label{Complexity.comp_copy_run}
\leanok
\uses{Complexity.CompReach, Complexity.LogProg.registerCopy_loop, Complexity.compARM, Complexity.compCopy, Complexity.compOracle}
\difficulty{3}
Let $C,e,B$ be natural numbers, let $m$ be a Boolean, let $f$ map Boolean words to
Boolean words, and let $h$ be a Boolean word. Choose a stage $s$ among eight stages and
let $v$ assign each of the fourteen registers a natural number at most $B$. Suppose the
source register specified by stage $s$ is at most the destination register specified by
that stage. Then the abstract machine, starting in the copy state for stage $s$ with its
bump flag false, register assignment $v$, and no recorded answer, reaches the state
prescribed by that stage with the destination register set to the original
source-register value and every other register unchanged; throughout the run, the
prescribed configuration invariant holds.
\end{lemma}

\begin{lemma}[Pair Parsing Run]
\lean{Complexity.comp_pair_run}
\label{Complexity.comp_pair_run}
\leanok
\uses{Complexity.CompReach, Complexity.LogProg.alignedPair_loop, Complexity.compARM, Complexity.compAdvance, Complexity.compLimit, Complexity.compOracle, Complexity.compStride, Complexity.comp_bit, Turing.pairDecode}
\difficulty{4}
Let $C,e$ be natural numbers, let $m$ and $i$ be Boolean parameters, let $f$ map Boolean
words to Boolean words, and let $h,w$ be Boolean words. Let $B$ be a natural number and
let $v$ assign a natural number to each of fourteen registers, with every register value
at most $B$. Suppose the register indexed by $1$ has value $v(0)+s|w|$, where $s=2$ if
$i$ is true and $s=1$ otherwise, and suppose that for every $j<|w|$, the bit of $h$ at
position $v(0)+sj$ equals the bit of $w$ at position $j$, interpreting out-of-range bits
as false. If the aligned pair parser does not parse $w$, then for some $t\le |w|$ there
is a bounded safe run from the configuration with state “scan” parameterized by $i$,
register values $v$, and no result, to one with state “answer false,” register $0$
increased by $st$, all other register values unchanged, and no result. If the parser
returns a pair $(a,b')$, there is a bounded safe run from that same initial
configuration to one with the copy state at stage $4$ when $i$ is true and stage $0$
otherwise, register $0$ increased by $s(2|a|+2)$, all other register values unchanged,
and no result.
\end{lemma}

\begin{lemma}[Halting from an Answer State]
\lean{Complexity.comp_ret}
\label{Complexity.comp_ret}
\leanok
\uses{Complexity.CompHalt, Complexity.LogProg.AHalt.ret, Complexity.comp_preS}
\difficulty{1}
Let $C,e\in\mathbb{N}$, let $\mathrm{mode}$ and $b$ be Boolean values, let $f$ map
finite Boolean strings to finite Boolean strings, let $h$ be a finite Boolean string,
and let $B\in\mathbb{N}$. For any vector $v$ of fourteen natural-number registers whose
entries are at most $B$, the computation started in the answer state for $b$, with
register vector $v$ and no result yet, halts with result $b$. Every configuration
strictly before halting satisfies the machine precondition and has all registers at most
$B$.
\end{lemma}

\begin{lemma}[Addressing the Second Component]
\lean{Complexity.comp_suffix_address}
\label{Complexity.comp_suffix_address}
\leanok
\uses{Complexity.pairEncode_getD, Complexity.pairEncode_getD_snd, Turing.length_pairEncode, Turing.pairEncode}
\difficulty{3}
For any bit strings $x,b,\mathrm{tail}$ and any natural number $j<|b|$, the bit at
position $4|x|+4+2j$ in the encoding of the encoding of $x$ paired with $b$, then paired
with $\mathrm{tail}$, is the bit of $b$ at position $j$.
\end{lemma}

\begin{lemma}[Last bit of a nested encoding]
\lean{Complexity.comp_last_address}
\label{Complexity.comp_last_address}
\leanok
\uses{Complexity.comp_suffix_address, Turing.length_pairEncode, Turing.pairEncode}
\difficulty{3}
Let $x$, $b$, and $t$ be bit strings, with $b$ nonempty. In the encoding of the pair
whose first component is the encoding of $(x,b)$ and whose second component is $t$, the
bit at position $2\bigl(\text{length of the encoding of }(x,b)\bigr)-2$, using false as the
default for an out-of-range position, equals the last bit of $b$, using true as the
default when taking the last bit.
\end{lemma}

\begin{lemma}[Canonical payload check run]
\lean{Complexity.comp_canonical_run}
\label{Complexity.comp_canonical_run}
\leanok
\uses{Complexity.CompReach, Complexity.LogProg.AReach.single, Complexity.LogProg.AReach.step, Complexity.LogProg.AReach.trans, Complexity.LogProg.Canon, Complexity.LogProg.PreS, Complexity.LogProg.astep, Complexity.LogProg.drop_canon, Complexity.LogProg.register_update_le, Complexity.LogProg.vword_whole₁, Complexity.compARM, Complexity.compCopy, Complexity.compOracle, Complexity.comp_bit, Complexity.comp_copy_run, Complexity.comp_last_address, Complexity.comp_preS, Turing.length_pairEncode, Turing.pairEncode}
\difficulty{5}
Let $C,e,B$ be natural numbers, let $m$ be a Boolean, let $f$ map Boolean words to
Boolean words, and let $x,b,$ and $t$ be Boolean words. Let $v$ assign each of the
fourteen registers a natural number at most $B$, with register $3$ equal to twice the
length of the encoding of $x$ paired with $b$, and register $4$ equal to $4|x|+4$. Then
there is a register assignment $v'$, also bounded by $B$, that agrees with $v$ on every
register other than register $0$, such that the abstract machine, starting in the
canonicality-check state on the input encoding of the pair whose first component is the
encoding of $x$ paired with $b$ and whose second component is $t$, reaches the
inner-preparation state if $b$ is canonical and otherwise the negative-answer state,
with no answer recorded; throughout the run, the prescribed configuration invariant
holds.
\end{lemma}

\begin{lemma}[Outer preparation run]
\lean{Complexity.comp_outer_prepare}
\label{Complexity.comp_outer_prepare}
\leanok
\uses{Complexity.CompReach, Complexity.LogProg.AReach.refl, Complexity.LogProg.AReach.step, Complexity.LogProg.AReach.trans, Complexity.LogProg.PreS, Complexity.LogProg.astep, Complexity.compARM, Complexity.compBound, Complexity.compCopy, Complexity.compFactor, Complexity.compOracle, Complexity.compValues, Complexity.comp_bound_room, Complexity.comp_copy_run, Complexity.comp_poly_run, Complexity.comp_power_bound, Complexity.comp_preS, Turing.length_pairEncode, Turing.pairEncode}
\difficulty{5}
Let $C,e$ be natural numbers, let $m$ be a Boolean, let $f$ map Boolean words to Boolean
words, and let $z,t$ be Boolean words. Set $h=\langle z,t\rangle$, $a=2|z|+2$,
$c=(2C+3)(|h|+1)^{e+1}$, and $B=8(2C+3)(|h|+1)^{e+1}$. Starting with register $0$ equal
to $a$, all other registers zero, no result, and state “copy” at stage $0$ with no
increment pending, there is a finite run to state “tail test” with its flag true and no
result. At the endpoint, registers $0,2,3,7,8$ equal $a,a,2|z|,c,|h|+1$, respectively,
register $1$ equals $|h|$, and all other registers are zero; every configuration before
the endpoint satisfies the program’s state precondition and has every register at most
$B$.
\end{lemma}

\begin{lemma}[Outer decoding run or rejection]
\lean{Complexity.comp_outer_decode}
\label{Complexity.comp_outer_decode}
\leanok
\uses{Complexity.CompHalt, Complexity.CompReach, Complexity.LogProg.AReach.halt, Complexity.LogProg.AReach.refl, Complexity.LogProg.AReach.single, Complexity.LogProg.AReach.step, Complexity.LogProg.AReach.trans, Complexity.LogProg.Canon, Complexity.LogProg.astep, Complexity.LogProg.bitsVal, Complexity.LogProg.boundedDecode_loop, Complexity.LogProg.canonicalTail_loop, Complexity.LogProg.drop_canon, Complexity.LogProg.register_update_le, Complexity.compARM, Complexity.compBound, Complexity.compCopy, Complexity.compDecodeAdvance, Complexity.compOracle, Complexity.compValues, Complexity.comp_bit, Complexity.comp_bound_room, Complexity.comp_copy_run, Complexity.comp_preS, Complexity.comp_ret, Complexity.pairEncode_getD, Complexity.query_accumulator_bound, Turing.length_pairEncode, Turing.pairEncode}
\difficulty{6}
Let $C,e$ be natural numbers, let $m$ be a Boolean, let $f$ map finite Boolean words to
finite Boolean words, and let $z,t$ be finite Boolean words. Write $h=\langle z,t\rangle$, where the
pairing doubles each bit of $z$, appends the separator $01$, and then appends $t$; set
$a=2|z|+2$, $c=(2C+3)(|h|+1)^{e+1}$, and $B=8(2C+3)(|h|+1)^{e+1}$. Start with input $h$,
state tail test with its flag true, no result, and register values
$a,|h|,a,2|z|,0,0,0,c,|h|+1,0,0,0,0,0$ in registers $0$ through $13$, respectively. If
$t$ is canonical and $|t|$ is at most the length of the binary representation of $c$,
there is a finite run to state scan with its inner flag true and no result; every
configuration before the endpoint satisfies the program's state precondition and has all
register values at most $B$, and at the endpoint the register values are
$0,|h|,a,2|z|,0,\mathrm{bitsVal}(t),0,0,|h|+1,0,0,0,0,0$. Otherwise, there is a finite
run to a halted configuration with result false, every configuration before the endpoint
satisfying that same state precondition and register bound.
\end{lemma}

\begin{lemma}[Inner preparation run]
\lean{Complexity.comp_inner_prepare}
\label{Complexity.comp_inner_prepare}
\leanok
\uses{Complexity.CompReach, Complexity.LogProg.AReach.step, Complexity.LogProg.AReach.trans, Complexity.LogProg.astep, Complexity.LogProg.register_update_le, Complexity.compARM, Complexity.compBound, Complexity.compCopy, Complexity.compFactor, Complexity.comp_copy_run, Complexity.comp_poly_run, Complexity.comp_power_bound, Complexity.comp_preS}
\difficulty{4}
Let $C,e,n$ be natural numbers, let the mode be Boolean, let $f$ map finite Boolean
words to finite Boolean words, let $h$ be a finite Boolean word, and let $v$ assign
natural numbers to fourteen registers. Suppose every register value is at most
$8(2C+3)(|h|+1)^{e+1}$, register $4$ equals $4n+4$, registers $6,7,10,11$ are zero, and
$n\le |h|$. Starting in the inner-preparation state with no result, the computation
reaches the length-bound state with registers $8$ and $7$ equal to $n+1$ and $C(n+1)^e$,
respectively, and every other register unchanged. Every configuration before the final
one satisfies the prescribed invariant: the instruction is admissible at that
configuration and all register values are at most $8(2C+3)(|h|+1)^{e+1}$.
\end{lemma}

\begin{definition}[Complexity walk state]
\lean{Complexity.compWalk}
\label{Complexity.compWalk}
\leanok
\uses{Complexity.CompState}
For natural numbers $C,e$, a Boolean $\mathit{suffix}$, and a natural number $i$, define
a state as follows: if $i=0$, it is the walk state with flags $\mathit{suffix}$ and
false; otherwise, if $\mathit{suffix}$ is true and $i=1$, it is the walk state with both
flags true; in every other case, it is the seek state with flag $\mathit{suffix}$.
\end{definition}

\begin{lemma}[Search Run Register Bounds]
\lean{Complexity.comp_seek_run}
\label{Complexity.comp_seek_run}
\leanok
\uses{Complexity.CompReach, Complexity.LogProg.indexedSeek_loop, Complexity.compARM, Complexity.compOracle, Complexity.compStride, Complexity.compWalk}
\difficulty{4}
Let $C,e,B,\text{count}$ be natural numbers, let $\text{mode}$ and $\text{suffix}$ be
Boolean flags, let $f$ map finite Boolean words to finite Boolean words, and let $h$ be
a finite Boolean word. Let $v$ assign a natural number to each of $14$ registers.
Suppose every register value is at most $B$, the value of register $3$ when
$\text{suffix}$ is true (or register $11$ when it is false) is
$v_0+\text{compStride}(\text{suffix})\text{count}$, register $13$ plus $\text{count}$ is at
most $B$, and $v_{13}\le v_5$. If $v_5<v_{13}+\text{count}$, there is a run from the
search state with flag $\text{suffix}$ to the suffix-hit state when $\text{suffix}$ is
true and the prefix-hit state otherwise; at its endpoint, registers $0$ and $13$ equal
$v_0+\text{compStride}(\text{suffix})(v_5-v_{13})$ and $v_5$, respectively. If
$v_5\ge v_{13}+\text{count}$, there is a run to the false-answer state when
$\text{suffix}$ is true and the separator state with flag false otherwise; registers $0$
and $13$ then equal $v_3$ (or $v_{11}$, respectively) and $v_{13}+\text{count}$. In
either case, every configuration before the endpoint satisfies the program's
step-validity condition and has all register values at most $B$.
\end{lemma}

\begin{lemma}[Oracle call outcome]
\lean{Complexity.comp_call_return}
\label{Complexity.comp_call_return}
\leanok
\uses{Complexity.CompHalt, Complexity.CompState, Complexity.LogProg.AHalt.step, Complexity.LogProg.astep, Complexity.LogProg.vword_whole₁, Complexity.compARM, Complexity.compOracle, Complexity.comp_preS, Complexity.comp_ret, Turing.pairEncode}
\difficulty{3}
Let $C,e,B$ be natural numbers, let $m$ be a Boolean, let $f$ map Boolean words to
Boolean words, and let $h$ be a Boolean word. Let $v$ assign a natural number at most
$B$ to each of the fourteen registers, and let $l$ be a computation state, with oracle
index $j$ among four indices and register $r$ among the fourteen registers. Suppose the
instruction at $l$ calls oracle $j$ in whole-input mode with the single argument
register $r$, branching to the positive-answer state on true and the negative-answer
state on false. Then the abstract machine halts from state $l$, register assignment $v$,
and no recorded answer, with result equal to the oracle's answer on the encoding of $h$
paired with the binary word for $v(r)$; throughout the run, the prescribed configuration
invariant holds.
\end{lemma}

\begin{lemma}[Suffix Hit Halting Result]
\lean{Complexity.comp_suffix_hit}
\label{Complexity.comp_suffix_hit}
\leanok
\uses{Complexity.CompHalt, Complexity.LogProg.AHalt.ret, Complexity.comp_bit, Complexity.comp_call_return, Complexity.comp_preS}
\difficulty{3}
Let $C,e,B$ be natural numbers, let $\text{mode}$ be a Boolean flag, let $f$ map finite
Boolean words to finite Boolean words, let $h$ be a finite Boolean word, and let $v$
assign a natural number to each of $14$ registers. Suppose every register value is at
most $B$. Starting in the suffix-hit state with register vector $v$ and no answer
recorded, the machine halts with answer true if $\text{mode}$ is true, and otherwise with
the bit of $h$ at position $v_0$, using false when that position is beyond the end;
every configuration before halting satisfies the program's step-validity condition and
has all register values at most $B$.
\end{lemma}

\begin{lemma}[Prefix-match halt result]
\lean{Complexity.comp_prefix_hit}
\label{Complexity.comp_prefix_hit}
\leanok
\uses{Complexity.CompHalt, Complexity.LogProg.AHalt.ret, Complexity.LogProg.AHalt.step, Complexity.LogProg.astep, Complexity.LogProg.register_update_le, Complexity.compARM, Complexity.comp_call_return, Complexity.comp_lifted, Complexity.comp_preS, Turing.pairEncode}
\difficulty{4}
Let $C,e\in\mathbb{N}$, let $m$ be a Boolean value, let $f$ map finite bit strings to
finite bit strings, let $x,b$ be finite bit strings, let $p,B\in\mathbb{N}$, and let $v$
be a vector of fourteen natural-number registers whose values are all at most $B$. On
the input encoding $\langle\langle x,b\rangle,\mathrm{bits}(p)\rangle$, the specified
machine started in the prefix-hit state with register vector $v$ halts with answer true
if $m$ is true, and otherwise with answer equal to the bit of $f(x)$ at index $\lfloor
v(0)/2\rfloor$, taking an out-of-range index to have value false. Every configuration
before halting satisfies the machine's admissibility condition and has all register
values at most $B$.
\end{lemma}

\begin{lemma}[Suffix Search Computation]
\lean{Complexity.comp_suffix_run}
\label{Complexity.comp_suffix_run}
\leanok
\uses{Complexity.CompHalt, Complexity.CompReach, Complexity.LogProg.AReach.halt, Complexity.LogProg.register_update_le, Complexity.compStride, Complexity.comp_ret, Complexity.comp_seek_run, Complexity.comp_suffix_address, Complexity.comp_suffix_hit, Turing.length_pairEncode, Turing.pairEncode}
\difficulty{5}
Let $C,e$ be natural numbers, let $\mathrm{mode}$ be a Boolean, let $f$ map finite
Boolean words to finite Boolean words, and let $x,b$ be finite Boolean words. Let $p,B$
be natural numbers and let $v$ assign a natural number to each of fourteen registers,
with every register value at most $B$. Suppose register $0$ contains $4|x|+4$, register
$3$ contains twice the length of $\langle x,b\rangle$, register $13$ contains
$2|f(x)|+2$, and register $5$ contains $p$. Assume also that $2|f(x)|+2\le p$ and
$2|f(x)|+2+|b|\le B$. Starting in the suffix-search state with register values $v$, no
result yet, and input $\langle\langle x,b\rangle,\mathrm{bits}(p)\rangle$, the counter
machine halts with result true precisely when $p<2|f(x)|+2+|b|$ if $\mathrm{mode}$ is
true; if it is false, the result is the bit of $b$ at index $p-(2|f(x)|+2)$, or false
when that index is beyond the end of $b$. Before halting, every configuration satisfies
the machine's step-validity condition, and every register remains at most $B$.
\end{lemma}

\begin{lemma}[Separator Search Halting Result]
\lean{Complexity.comp_separator_run}
\label{Complexity.comp_separator_run}
\leanok
\uses{Complexity.CompHalt, Complexity.CompReach, Complexity.LogProg.AHalt.step, Complexity.LogProg.AReach.halt, Complexity.LogProg.AReach.single, Complexity.LogProg.AReach.step, Complexity.LogProg.AReach.trans, Complexity.LogProg.astep, Complexity.LogProg.register_update_le, Complexity.compARM, Complexity.compCopy, Complexity.compOracle, Complexity.comp_copy_run, Complexity.comp_preS, Complexity.comp_ret, Complexity.comp_suffix_run, Complexity.pairEncode_getD, Turing.length_pairEncode, Turing.pairEncode}
\difficulty{6}
Let $C,e,B,p$ be natural numbers, let $\text{mode}$ be a Boolean flag, let $f$ map finite
Boolean words to finite Boolean words, and let $x,b$ be finite Boolean words. Let $v$
assign a natural number to each of $14$ registers. Suppose every register value is at
most $B$, register $3$ equals twice the length of the pairing of $x$ with $b$, register
$4$ equals $4|x|+4$, register $13$ equals twice the length of $f(x)$, register $5$
equals $p$, twice the length of $f(x)$ is at most $p$, and the length of the pairing of
$f(x)$ with $b$ is at most $B$. Starting in the separator state with flag false and with
no answer recorded, the machine halts with answer $p<|\langle f(x),b\rangle|$ if
$\text{mode}$ is true, and otherwise with the bit at position $p$ of $\langle
f(x),b\rangle$, taking the bit to be false when $p$ is beyond its length. Every
configuration before halting satisfies the program's step-validity condition and has all
register values at most $B$.
\end{lemma}

\begin{lemma}[Encoded output lookup run]
\lean{Complexity.comp_lookup_run}
\label{Complexity.comp_lookup_run}
\leanok
\uses{Complexity.CompHalt, Complexity.CompReach, Complexity.LogProg.AReach.halt, Complexity.LogProg.register_update_le, Complexity.compStride, Complexity.comp_prefix_hit, Complexity.comp_seek_run, Complexity.comp_separator_run, Complexity.pairEncode_getD, Turing.length_pairEncode, Turing.pairEncode}
\difficulty{6}
Let $C,e,p,B$ be natural numbers, let $m$ be a Boolean mode, let $f$ map finite bit
strings to finite bit strings, let $x,b$ be finite bit strings, and let $v$ assign a
natural number to each of fourteen registers. Suppose every register value is at most
$B$, register $0$ is $0$, register $3$ is twice the length of the encoding of $(x,b)$,
register $4$ is $4|x|+4$, register $11$ is twice the length of $f(x)$, register $13$ is
$0$, and register $5$ is $p$. Suppose also that the encoding of $(f(x),b)$ has length at
most $B$. Then, from the seek state on the nested encoding of $(x,b)$ with binary index
$p$, register vector $v$, and no answer, the computation halts with answer equal to the
Boolean that is true exactly when $p$ is less than the length of the encoding of
$(f(x),b)$ if $m$ is true, and otherwise equal to the bit at position $p$ of that
encoding, with out-of-range positions interpreted as false. Every configuration before
halting satisfies the program’s instruction-side conditions and register bound $B$.
\end{lemma}

\begin{lemma}[The query on an encoded pair]
\lean{Complexity.compQ_field}
\label{Complexity.compQ_field}
\leanok
\uses{Complexity.LogProg.Canon, Complexity.LogProg.bitsVal, Complexity.LogProg.canon_bits, Complexity.LogProg.canon_eq_bits, Complexity.compQ, Turing.pairDecode_pairEncode, Turing.pairEncode}
\difficulty{3}
Let $f$ map finite Boolean strings to finite Boolean strings, and let $x$ and $b$ be
finite Boolean strings. Applying the query transformation to the encoding of $(x,b)$
returns the encoding of $(f(x),b)$ if $b$ is empty or its last bit is true, and returns
the empty string otherwise.
\end{lemma}

\begin{definition}[Language from a Boolean function]
\lean{Complexity.compLang}
\label{Complexity.compLang}
\leanok
\uses{Complexity.compQ, Complexity.indexLang}
Given a function $f$ from Boolean lists to Boolean lists and a Boolean mode, define a
language of Boolean lists by indexing each input $z$ at position $p$: in true mode,
accept exactly when $p$ is less than the length of the query associated to $f$ and $z$;
in false mode, accept exactly when the bit at position $p$ in that query, using false
when the position is out of range, is true.
\end{definition}

\begin{lemma}[Membership of encoded pairs]
\lean{Complexity.comp_lang_pair}
\label{Complexity.comp_lang_pair}
\leanok
\uses{Complexity.compLang, Complexity.compQ, Complexity.indicator_indexLang_pairEncode, Turing.MultiTapeTM.indicator, Turing.pairEncode}
\difficulty{3}
Let $f$ be a function on finite bit strings, let $z$ be a bit string, let $p$ be a
natural number, and let $m$ be a Boolean mode. The Boolean indicator of the language
defined by testing the output associated to $z$ at index $p$, on the encoded input
consisting of $z$ and the binary representation of $p$, equals true exactly when $p$ is
less than the associated output's length if $m$ is true; when $m$ is false, it equals
the bit at position $p$ of that output, with out-of-range positions interpreted as
false.
\end{lemma}

\begin{lemma}[Indicator vanishes on malformed queries]
\lean{Complexity.comp_bad_query}
\label{Complexity.comp_bad_query}
\leanok
\uses{Complexity.LogProg.indicator_eq_decide, Complexity.compLang, Turing.MultiTapeTM.indicator, Turing.pairEncode}
\difficulty{4}
Let $f$ map bit strings to bit strings, let $h$ be a bit string, and suppose $h$ is not
the encoding of $(z,\operatorname{bin}(p))$ for any bit string $z$ and natural number
$p$. For either Boolean choice of mode, the indicator of the language that contains
encoded queries exactly when either $p$ is less than the transformed output length (in
one mode) or the output bit at position $p$ is true (in the other mode) evaluates to
false at $h$.
\end{lemma}

\begin{lemma}[Rejection of an Invalid Tail]
\lean{Complexity.comp_rejected_tail}
\label{Complexity.comp_rejected_tail}
\leanok
\uses{Complexity.LogProg.Canon, Complexity.LogProg.bitsVal, Complexity.LogProg.canon_bits, Complexity.LogProg.canon_eq_bits, Complexity.compLang, Complexity.compQ_query_overflow, Complexity.comp_bad_query, Complexity.comp_lang_pair, Turing.MultiTapeTM.indicator, Turing.pairEncode, Turing.pairEncode_injective}
\difficulty{5}
Let $f$ map finite bit strings to finite bit strings, and suppose there are natural
numbers $C,e$ such that $|f(x)|\le C(|x|+1)^e$ for every bit string $x$. For either
Boolean value $\mathrm{mode}$ and any bit strings $z$ and $\mathrm{tail}$, if it is not
the case that $\mathrm{tail}$ is canonical and its length is at most the length of the
binary representation of $(2C+3)(|\langle z,\mathrm{tail}\rangle|+1)^{e+1}$, then the
indicator of the specified language is false on $\langle z,\mathrm{tail}\rangle$.
\end{lemma}

\begin{lemma}[Inner Preparation Computation]
\lean{Complexity.comp_valid_inner}
\label{Complexity.comp_valid_inner}
\leanok
\uses{Complexity.CompHalt, Complexity.CompReach, Complexity.LogProg.AHalt.step, Complexity.LogProg.AReach.halt, Complexity.LogProg.AReach.trans, Complexity.LogProg.Canon, Complexity.LogProg.registerScale_loop, Complexity.LogProg.register_update_le, Complexity.compARM, Complexity.compBound, Complexity.compFactor, Complexity.compOracle, Complexity.compQ_field, Complexity.compQ_host_cap, Complexity.comp_bound_room, Complexity.comp_inner_prepare, Complexity.comp_length_search, Complexity.comp_lifted, Complexity.comp_lookup_run, Complexity.comp_power_bound, Complexity.comp_preS, Turing.length_pairEncode, Turing.pairEncode}
\difficulty{6}
Let $C,e$ be natural numbers, let $\mathrm{mode}$ be Boolean, and let $f$ map finite
Boolean words to finite Boolean words. Suppose that for every Boolean word $y$,
$|f(y)|\le C(|y|+1)^e$. Let $x,b$ be Boolean words, let $p$ be a natural number, and
suppose that $b$ is empty or its last bit is true. Let $v$ assign natural numbers to
fourteen registers, each at most $8(2C+3)(|\langle\langle
x,b\rangle,\mathrm{bits}(p)\rangle|+1)^{e+1}$. Suppose register $3$ equals $2|\langle
x,b\rangle|$, register $4$ equals $4|x|+4$, register $5$ equals $p$, and registers
$6,7,10,11,12,13$ are all zero. Starting in the inner-preparation state with these
register values, no result yet, and input $\langle\langle
x,b\rangle,\mathrm{bits}(p)\rangle$, the counter machine halts with result true
precisely when $p<|\langle f(x),b\rangle|$ if $\mathrm{mode}$ is true; otherwise its
result is the bit at index $p$ of $\langle f(x),b\rangle$, or false if that index is
beyond the end. Before halting, every configuration satisfies the machine's
step-validity condition, and all register values remain within the stated bound.
\end{lemma}

\begin{lemma}[Bounded inner computation]
\lean{Complexity.comp_inner_run}
\label{Complexity.comp_inner_run}
\leanok
\uses{Complexity.CompHalt, Complexity.CompReach, Complexity.LogProg.AReach.halt, Complexity.LogProg.AReach.trans, Complexity.LogProg.Canon, Complexity.LogProg.register_update_le, Complexity.compBound, Complexity.compCopy, Complexity.compLimit, Complexity.compQ, Complexity.compQ_field, Complexity.compStride, Complexity.compValues, Complexity.comp_bound_room, Complexity.comp_canonical_run, Complexity.comp_copy_run, Complexity.comp_pair_run, Complexity.comp_ret, Complexity.comp_valid_inner, Complexity.drop_first_read, Turing.eq_pairEncode_of_pairDecode, Turing.length_pairEncode, Turing.pairDecode, Turing.pairDecode_pairEncode, Turing.pairEncode}
\difficulty{6}
Let $C,e,p\in\mathbb N$, let $m$ be a Boolean mode, and let $f$ map finite Boolean words
to finite Boolean words with $|f(x)|\le C(|x|+1)^e$ for every such word $x$. For a
finite Boolean word $z$, suppose $p\le 8(2C+3)(|\langle
z,\mathrm{bits}(p)\rangle|+1)^{e+1}$, where $\mathrm{bits}(p)$ is the binary word for
$p$ and $\langle\cdot,\cdot\rangle$ is the self-delimiting pairing. The abstract
computation on input $\langle z,\mathrm{bits}(p)\rangle$, starting in the inner scan
state with registers $1$ through $14$ initialized respectively to $|\langle
z,\mathrm{bits}(p)\rangle|$, $2|z|+2$, $2|z|$, $0$, $p$, $0$, $0$, $|\langle
z,\mathrm{bits}(p)\rangle|+1$, $0$, $0$, $0$, $0$, $0$, and with no result recorded,
halts with result $\mathbf{1}_{\{p<|q|\}}$ when $m$ is true and with the bit at position
$p$ of $q$ (false if that position is absent) when $m$ is false, where $q$ is the word
obtained by decoding $z$ as $(x,b)$ and returning $\langle f(x),b\rangle$ if $b$ is the
canonical binary word for its decoded value, and the empty word otherwise. At every time
before it halts, all register values are at most $8(2C+3)(|\langle
z,\mathrm{bits}(p)\rangle|+1)^{e+1}$; moreover, each active instruction comparing two
registers compares distinct registers, instructions comparing a register with the input
word are used only on inputs of the required valid encoding (a unary prefix followed by
a canonical payload, or by a pair of canonical payloads, as appropriate), and every
oracle call uses a list of distinct registers.
\end{lemma}

\begin{lemma}[Bounded Computation of the Encoded Language]
\lean{Complexity.comp_run}
\label{Complexity.comp_run}
\leanok
\uses{Complexity.CompHalt, Complexity.CompReach, Complexity.LogProg.AReach.halt, Complexity.LogProg.AReach.trans, Complexity.LogProg.Canon, Complexity.LogProg.bitsVal, Complexity.LogProg.bitsVal_bits, Complexity.LogProg.canon_eq_bits, Complexity.LogProg.inputLength_loop, Complexity.LogProg.register_update_le, Complexity.compARM, Complexity.compBound, Complexity.compLang, Complexity.compLimit, Complexity.compOracle, Complexity.compStride, Complexity.compValues, Complexity.comp_bad_query, Complexity.comp_bound_room, Complexity.comp_inner_run, Complexity.comp_lang_pair, Complexity.comp_len, Complexity.comp_outer_decode, Complexity.comp_outer_prepare, Complexity.comp_pair_run, Complexity.comp_rejected_tail, Complexity.comp_ret, Complexity.query_index_bound, Turing.MultiTapeTM.indicator, Turing.eq_pairEncode_of_pairDecode, Turing.pairDecode, Turing.pairDecode_pairEncode, Turing.pairEncode}
\difficulty{6}
Let $C,e\in\mathbb{N}$, let $\mathrm{mode}$ be Boolean, and let $f$ map finite Boolean
strings to finite Boolean strings. Suppose that for every finite Boolean string $x$,
$|f(x)|\le C(|x|+1)^e$. For every finite Boolean string $h$, the computation started in
its counting state with all fourteen registers zero and no result yet halts with the
indicator of the following language on $h$: strings encoding a pair $(z,p)$, where $p$
is the binary expansion of a natural number, such that, on decoding $z$ as a pair
$(x,b)$, the binary expansion of the natural number represented by $b$ equals $b$, and
either $p$ is less than the length of the encoding of $(f(x),b)$ when $\mathrm{mode}$ is
true, or the bit of that encoding at index $p$ is true when $\mathrm{mode}$ is false.
Throughout the run, every configuration strictly before halting satisfies the machine
precondition and has all registers at most $8(2C+3)(|h|+1)^{e+1}$.
\end{lemma}

\begin{lemma}[Logspace membership of the constructed language]
\lean{Complexity.comp_mem}
\label{Complexity.comp_mem}
\leanok
\uses{Complexity.ImplicitlyLogspaceComputable, Complexity.LOGSPACE, Complexity.LogProg.arm_decides_poly, Complexity.compARM, Complexity.compBank, Complexity.compBank_mem, Complexity.compLang, Complexity.comp_run}
\difficulty{3}
Let $f$ be a function on bit strings such that its output length is bounded by a
polynomial in the input length, the language of pairs $(x,i)$ for which the $i$th bit of
$f(x)$ is true is decidable in logarithmic space, and the language of pairs $(x,i)$ for
which $i<|f(x)|$ is decidable in logarithmic space. Let $C,e$ be natural numbers such
that $|f(x)|\le C(|x|+1)^e$ for every bit string $x$, and let $m$ be a Boolean mode.
Then the language that, on an encoded pair $(z,p)$, tests whether $p$ is below the
length of the output associated to $z$ when $m$ is true, and otherwise tests the output
bit at position $p$, is decidable in logarithmic space.
\end{lemma}

\begin{theorem}[Implicit logspace closure under guarded composition]
\lean{Complexity.ImplicitlyLogspaceComputable.compQ}
\label{Complexity.ImplicitlyLogspaceComputable.compQ}
\leanok
\uses{Complexity.ImplicitlyLogspaceComputable, Complexity.compLang, Complexity.compQ, Complexity.compQ_length_le, Complexity.comp_mem}
\difficulty{3}
Let $f$ map Boolean words to Boolean words and suppose its output length is bounded by a
polynomial in the input length, while the languages of encoded pairs $\langle
x,\operatorname{bits}(i)\rangle$ for which the $i$-th bit of $f(x)$ is true and for
which $i<|f(x)|$ are decidable in logarithmic space. Define a function on Boolean words
$w$ by parsing $w$ as $\langle x,b\rangle$; if parsing succeeds and $b$ is the binary
representation of its own value, return $\langle f(x),b\rangle$, and otherwise return
the empty word. This function has polynomially bounded output length, and its true-bit
index language and output-length index language are both decidable in logarithmic space.
\end{theorem}
